% ============================================================================
%  DDS-TS 1.0 — ZeroDDS Vendor Specification.
%
%  TypeScript Platform-Specific Model for OMG IDL 4.2 + X-Types 1.3.
%
%  Build:  tectonic main.tex
%  Output: main.pdf  (renamed to dds-ts-1.0-beta1.pdf in CI release jobs)
% ============================================================================

\documentclass{../styles/zerodds-vendor-spec}

\title{DDS-TS\\\large TypeScript Platform-Specific Model for OMG IDL 4.2}
\specshorttitle{DDS-TS 1.0}
\specversion{zerodds/v1.0}
\specstatus{Final}
\specdate{2026-04-29}
\speceditors{Sandra Ke{\ss}ler}

\begin{document}

\maketitle

\tableofcontents
\clearpage

% ============================================================================
%  Status of this Document.
% ============================================================================
\chapter*{Status of this Document}
\addcontentsline{toc}{chapter}{Status of this Document}

This specification (version \texttt{zerodds/v1.0}) is
published by the ZeroDDS Project. It defines a Platform-Specific
Model (PSM) that maps OMG IDL 4.2 to TypeScript. Its document
organisation follows the OMG language PSMs (IDL-to-C++,
IDL-to-Java, IDL-to-C\#); its type-system mapping deliberately
diverges, using TypeScript's structural interfaces rather than
the nominal classes of those PSMs
(\S\,\ref{sec:mapping-philosophy}).

This document is not an OMG document and is not on any OMG
standardisation track. The conformance surface is nevertheless
designed for cross-vendor interoperability: the brand schemas
(\S\,\ref{sec:wasm-types}, \S\,\ref{sec:char-brand}), exported
symbol names, operation/exception encodings
(Chapter~\ref{ch:operations}), and diagnostic codes
(Annex~\ref{annex:diagnostics}) are vendor-neutral.

TypeScript is not an OMG or ISO standard; its prose
specification (Microsoft, last edition 1.8 / 2016) is no longer
maintained. The de-facto reference is the \texttt{tsc} compiler
distributed as the \texttt{typescript} npm package. A conformant
generator output \rfckw{SHALL} compile cleanly under
\texttt{tsc --strict} for any \texttt{typescript} package
release at major version 5 (i.e.\ \texttt{5.0.0} up to but not
including \texttt{6.0.0}). Runtime semantics are inherited from ECMA-262
(14th edition or later).

The terms \rfckw{MUST}, \rfckw{MUST NOT}, \rfckw{SHALL},
\rfckw{SHALL NOT}, \rfckw{SHOULD}, \rfckw{SHOULD NOT}, \rfckw{MAY},
\rfckw{MAY NOT}, \rfckw{REQUIRED}, \rfckw{RECOMMENDED}, and
\rfckw{OPTIONAL} are to be interpreted as described in IETF RFC
2119 and RFC 8174.

% ============================================================================
%  Chapter 1 — Scope.
% ============================================================================
\chapter{Scope}

\section{Purpose}

This specification defines a TypeScript Platform-Specific Model
(PSM) for the OMG Interface Definition Language version 4.2, with
extensions for the Extensible and Dynamic Topic Types for DDS
(X-Types) version 1.3 annotations.

The mapping enables IDL types defined for DDS topics to be
generated as TypeScript declarations and consumed in browser-side
or Node.js-side TypeScript code without loss of semantic
information.

\section{Out of Scope}

The following items are explicitly outside the scope of this
specification:

\begin{itemize}
  \item TypeScript runtime support beyond the descriptor-runtime
        library defined in Chapter~\ref{ch:descriptor-runtime}.
  \item Redefinition of any DDS wire-encoding format. Wire
        encoding is governed by OMG X-Types (XCDR1, XCDR2); this
        document references those rules but does not redefine
        them.
  \item Bindings to a specific JavaScript runtime: this
        specification produces ESM source compilable under
        \texttt{tsc --strict}.
  \item DDS-RPC wire protocol; covered by OMG DDS-RPC. This PSM
        only specifies the TypeScript surface through which RPC
        client and server code is generated
        (Chapter~\ref{ch:operations}).
\end{itemize}

The mapping covers the full structural surface of OMG IDL 4.2:
modules, free-standing constants, primitive types, constructed
types (struct, union, enum, bitset, bitmask, sequence, array,
typedef, exception), and the \texttt{interface} construct with
its operations and attributes.

% ============================================================================
%  Chapter 2 — Conformance.
% ============================================================================
\chapter{Conformance}\label{ch:conformance}

This specification defines four conformance profiles. A
conformance claim \rfckw{MUST} cite this specification by full
identifier (\texttt{DDS-TS, zerodds/v1.0} or later) and
enumerate every profile to which conformance is claimed.

\section{IDL-Codegen Profile}\label{sec:profile-codegen}

\begin{conformance}{IDL-Codegen Profile}
An implementation conforms to the IDL-Codegen Profile if and only
if all of the following are true:
\begin{enumerate}
  \item It implements the type mapping defined in
        Chapter~\ref{ch:type-mapping} for every primitive,
        constructed, and template type of OMG IDL 4.2.
  \item It implements the annotation mapping defined in
        \S\,\ref{sec:annotation-mapping} for every X-Types 1.3
        annotation listed in that section.
  \item For every IDL type it emits, it additionally emits a
        \emph{Type Descriptor} (Chapter~\ref{ch:descriptor-runtime})
        and a corresponding type-guard function, per
        \S\,\ref{sec:struct-mapping}.
  \item It produces output that is valid TypeScript $\geq$ 5.0
        source code, compilable with \texttt{tsc} \texttt{--strict}
        without the \texttt{experimentalDecorators} or
        \texttt{emitDecoratorMetadata} compiler options.
  \item It satisfies the conformance test cases in Annex~B,
        \S\,B.1.
\end{enumerate}
\end{conformance}

\section{Descriptor-Runtime Profile}\label{sec:profile-runtime}

\begin{conformance}{Descriptor-Runtime Profile}
An implementation conforms to the Descriptor-Runtime Profile if
and only if all of the following are true:
\begin{enumerate}
  \item It provides the descriptor types
        \texttt{DdsTypeDescriptor}, \texttt{DdsMemberDescriptor},
        \texttt{DdsTypeRef}, \texttt{ExtensibilityKind},
        \texttt{DdsAny}, with the structural shapes defined in
        Chapter~\ref{ch:descriptor-runtime}.
  \item It provides the reflection functions
        \texttt{getKey}, \texttt{getTopic}, \texttt{registerType},
        \texttt{lookupType}, \texttt{unboxAny}, \texttt{boxAny}
        with the semantics defined in
        Chapter~\ref{ch:descriptor-runtime}.
  \item It provides the runtime helpers \texttt{makeChar} and
        \texttt{makeWChar} per \S\,\ref{sec:char-brand}.
  \item It satisfies the conformance test cases in Annex~B,
        \S\,B.2.
\end{enumerate}
\end{conformance}

\section{Operations Profile}\label{sec:profile-operations}

\begin{conformance}{Operations Profile}
An implementation conforms to the Operations Profile if and only
if all of the following are true:
\begin{enumerate}
  \item It conforms to the IDL-Codegen Profile
        (\S\,\ref{sec:profile-codegen}) and to the
        Descriptor-Runtime Profile
        (\S\,\ref{sec:profile-runtime}).
  \item It implements the \texttt{interface}, operation,
        attribute, and exception mapping defined in
        Chapter~\ref{ch:operations}.
  \item For every IDL \texttt{interface} it emits, it additionally
        emits a corresponding \emph{Service Descriptor}
        (\S\,\ref{sec:operations-descriptor}).
  \item It provides the operations descriptor types
        \texttt{ServiceDescriptor}, \texttt{OperationDescriptor},
        \texttt{OperationParameterDescriptor},
        \texttt{AttributeDescriptor}, \texttt{ParameterMode},
        and the runtime helper \texttt{isOneOf}, with the
        structural shapes defined in
        \S\,\ref{sec:operations-descriptor}.
  \item It satisfies the conformance test cases in Annex~B,
        \S\,B.3.
\end{enumerate}
The Operations Profile is independent of the WASM-Bindings
Profile; an implementation \rfckw{MAY} claim either, neither, or
both.
\end{conformance}

\section{WASM-Bindings Profile}\label{sec:profile-wasm}

\begin{conformance}{WASM-Bindings Profile}
An implementation conforms to the WASM-Bindings Profile if and
only if all of the following are true:
\begin{enumerate}
  \item It conforms to the IDL-Codegen Profile
        (\S\,\ref{sec:profile-codegen}) and to the
        Descriptor-Runtime Profile
        (\S\,\ref{sec:profile-runtime}).
  \item It exposes the DDS DCPS data-path operations enumerated in
        Annex~C, \S\,C.1, as TypeScript-callable functions backed
        by a WebAssembly module compiled from a DDS implementation
        (e.g.\ the ZeroDDS Rust core).
  \item It serialises and deserialises samples using the XCDR2
        wire format on both sides of the WASM boundary; the binding
        layer \rfckw{MUST NOT} expose JavaScript-native object
        representations as the canonical wire form.
  \item It satisfies the conformance test cases in Annex~C,
        \S\,C.2 (round-trip publish/subscribe through the WASM
        boundary).
\end{enumerate}
\end{conformance}

The design and operational semantics of the WASM-Bindings Profile
are normatively specified in Annex~C.

% ============================================================================
%  Chapter 3 — Normative References.
% ============================================================================
\chapter{Normative References}\label{sec:ts-reference}

\begin{longtable}{p{3.4cm}p{2.0cm}p{8.5cm}}
  \toprule
  \textbf{Identifier} & \textbf{Version} & \textbf{Title} \\
  \midrule
  \endhead
  OMG IDL              & 4.2         & Interface Definition Language \\
  OMG X-Types          & 1.3         & Extensible and Dynamic Topic Types for DDS \\
  OMG DDS              & 1.4         & Data Distribution Service for Real-Time Systems \\
  OMG DDS-RPC          & 1.0         & Remote Procedure Call over DDS \\
  OMG IDL4-Cpp         & 1.0         & IDL to C++ Mapping (OMG language PSM) \\
  OMG IDL4-Csharp      & 1.0         & IDL to C\# Mapping (OMG language PSM) \\
  OMG IDL4-Java        & 1.0         & IDL to Java Mapping (OMG language PSM) \\
  ECMA-262             & 14th ed.    & ECMAScript Language Specification (TC39) \\
  Microsoft \texttt{tsc} & 5.0+      & TypeScript compiler behaviour (\texttt{typescript} npm package, github.com/microsoft/TypeScript); see \S\,\ref{sec:ts-reference} \\
  IETF RFC 2119        & ---         & Key words for use in RFCs \\
  IETF RFC 8174        & ---         & Ambiguity of Uppercase vs Lowercase \\
  \bottomrule
\end{longtable}

% ============================================================================
%  Chapter 4 — Terms and Definitions.
% ============================================================================
\chapter{Terms and Definitions}

\begin{description}[leftmargin=2cm,style=nextline]
  \item[IDL] OMG Interface Definition Language as defined in
        \omgref{OMG IDL}{4.2}.
  \item[Codegen] An implementation conforming to the IDL-Codegen
        Profile (\S\,\ref{sec:profile-codegen}).
  \item[Type Descriptor] A side-table value of TypeScript type
        \texttt{DdsTypeDescriptor<T>} that captures the metadata
        of an IDL type (name, extensibility, members, topic
        binding, type-guard); see Chapter
        \ref{ch:descriptor-runtime}.
  \item[Service Descriptor] A side-table value of TypeScript type
        \texttt{ServiceDescriptor<TClient,\,THandler>} that
        captures the metadata of an IDL \texttt{interface} (name,
        operations, attributes, exception sets, inheritance); see
        Chapter \ref{ch:operations}.
  \item[ESM] ECMAScript Module syntax: \texttt{export} and
        \texttt{import} statements as defined in ECMA-262.
\end{description}

% ============================================================================
%  Chapter 5 — Symbols.
% ============================================================================
\chapter{Symbols and Abbreviations}

\begin{description}[leftmargin=2cm]
  \item[DDS] Data Distribution Service
  \item[ESM] ECMAScript Module
  \item[IDL] Interface Definition Language
  \item[PSM] Platform-Specific Model
  \item[TS] TypeScript
  \item[WASM] WebAssembly
  \item[XCDR2] Extended Common Data Representation, version 2
\end{description}

% ============================================================================
%  Chapter 6 — Overview.
% ============================================================================
\chapter{Overview (Informative)}

\section{Mapping Philosophy}\label{sec:mapping-philosophy}

The DDS-TS mapping uses TypeScript's structural typing rather
than the nominal class-based mappings of the OMG C++, Java, and
C\# PSMs. Every IDL data type is mapped to an
\texttt{interface}; annotation metadata is exposed via a
separately-emitted \emph{Type Descriptor}
(Chapter~\ref{ch:descriptor-runtime}). The same approach
applies to IDL \texttt{interface} mappings
(Chapter~\ref{ch:operations}), which produce paired
client/handler interfaces plus a \emph{Service Descriptor}. No
TypeScript class is emitted by this PSM in either path.

\section{Runtime Reflection Strategy}\label{sec:reflection-strategy}

TypeScript interfaces are erased during compilation; a plain
object literal carries no JavaScript-level type identity. The
PSM closes this gap with first-class side-table values:

\begin{itemize}[noitemsep]
  \item Every IDL type emits an
        \texttt{export const $\langle$Type$\rangle$Type:\,DdsTypeDescriptor<T>}
        constant, registered with a process-wide
        \texttt{DdsTypeRegistry}
        (\S\,\ref{sec:descriptor-reflection}).
  \item Every IDL interface emits an
        \texttt{export const $\langle$Iface$\rangle$Service:\,
        ServiceDescriptor<C,\,H>} constant.
  \item Every IDL \texttt{any} value crosses the TypeScript
        surface as a \texttt{DdsAny\,\{ typeId, value \}} carrier
        whose \texttt{typeId} indexes the registry.
\end{itemize}

This is a deliberate parallel to the runtime type system that
OMG X-Types 1.3 itself defines for cross-vendor type discovery:
the \texttt{TypeObject} / \texttt{MinimalTypeObject} structures
of OMG X-Types 1.3 §7.3 carry exactly the same metadata
(extensibility, member ids, key membership, type references) as
first-class values that DDS implementations exchange over the
Type-Lookup service. The DDS-TS \texttt{DdsTypeDescriptor} is
the TypeScript-side projection of that runtime metadata: the
content set is intentionally close to \texttt{TypeObject}, so
that an implementation \rfckw{MAY} construct a
\texttt{TypeObject} from a \texttt{DdsTypeDescriptor} (and vice
versa) without information loss.

The side-table approach is necessary because TypeScript has no
Plain-Object-Reflection mechanism comparable to Java's
\texttt{Class<T>} or C\#'s \texttt{System.Type}: structural
interface-based code carries no runtime type identity, so the
identity must be carried alongside the data as a first-class
value. This PSM standardises the shape of that sidecar.

\section{Toolchain}

The reference implementation in \texttt{crates/idl-ts/} produces
\texttt{.ts} source files compatible with TypeScript $\geq$ 5.0. The
generated source uses ECMAScript-Module (ESM) syntax exclusively.

\section{Code-Style Toggle}

By default, IDL identifier names are preserved verbatim
(snake\_case to snake\_case, CamelCase to CamelCase). An optional
\emph{camelCase-toggle} configuration (Chapter~\ref{ch:code-style})
re-cases snake\_case IDL fields to camelCase TypeScript fields.

% ============================================================================
%  Chapter 7 — Type Mapping.
% ============================================================================
\chapter{Type Mapping (Normative)}\label{ch:type-mapping}

\section{Module Mapping}\label{sec:module-mapping}

An IDL \texttt{module} \rfckw{SHALL} be mapped to a TypeScript
\texttt{export namespace} declaration of the same name. Nested
\texttt{module} declarations \rfckw{SHALL} produce nested
\texttt{export namespace} declarations preserving the IDL nesting
order.

\begin{specexample}
\begin{lstlisting}
module Vehicle {
  module Sensor {
    struct Reading { double value; };
  };
  struct Car { string vin; };
};
\end{lstlisting}
maps to
\begin{lstlisting}
export namespace Vehicle {
  export namespace Sensor {
    export interface Reading { value: number; }
  }
  export interface Car { vin: string; }
}
\end{lstlisting}
\end{specexample}

A reference to a qualified IDL name (\texttt{::Vehicle::Sensor::Reading})
\rfckw{SHALL} be emitted as the corresponding dotted TypeScript path
(\texttt{Vehicle.Sensor.Reading}). Implementations \rfckw{MAY}
additionally produce one ECMAScript module file per top-level IDL
module under a \texttt{--module-per-file} configuration option; the
namespace structure inside each file \rfckw{SHALL} remain as
specified above.

A re-opening of the same module name across compilation units
\rfckw{SHALL} be honoured: TypeScript \texttt{export namespace}
permits declaration merging and the code generator \rfckw{SHALL}
rely on that mechanism rather than emitting a renamed namespace.

\section{Constant Mapping}\label{sec:constant-mapping}

A free-standing IDL constant declaration of the form
\texttt{const $\langle$type$\rangle$ $\langle$name$\rangle$\,=\,$\langle$expr$\rangle$;}
\rfckw{SHALL} be mapped to a TypeScript
\texttt{export const $\langle$name$\rangle$:\,T'\,=\,$\langle$expr'$\rangle$;}
declaration, where \texttt{T'} is the TypeScript-mapped form of
\texttt{$\langle$type$\rangle$} per
\S\,\ref{sec:primitive-mapping} and \texttt{$\langle$expr'$\rangle$}
is the const-evaluated literal in TypeScript syntax.

\begin{specexample}
\begin{lstlisting}
const long    MAX_RETRIES   = 5;
const double  PI            = 3.14159265358979;
const string  SERVICE_NAME  = "telemetry";
const boolean DEBUG_ENABLED = FALSE;
\end{lstlisting}
maps to
\begin{lstlisting}
export const MAX_RETRIES:   number  = 5;
export const PI:            number  = 3.14159265358979;
export const SERVICE_NAME:  string  = "telemetry";
export const DEBUG_ENABLED: boolean = false;
\end{lstlisting}
\end{specexample}

The implementation \rfckw{SHALL} const-evaluate the right-hand-side
expression at code-generation time. Decimal, hexadecimal, octal, and
binary integer-literal forms \rfckw{MUST} be supported, as well as
character escapes inside string literals (per IDL §7.2.6.5).
Constants of type \texttt{long long} or \texttt{unsigned long long}
\rfckw{SHALL} be emitted with a trailing \texttt{n} BigInt literal:
\texttt{export const X: bigint = 9007199254740993n;}.

A constant whose right-hand side references another constant
\rfckw{SHALL} be emitted using the referenced TypeScript identifier;
the implementation \rfckw{SHALL} resolve the dependency in
declaration order and \rfckw{SHALL} reject cyclic references at
parse time. Constants declared inside an IDL \texttt{module}
\rfckw{SHALL} appear inside the corresponding TypeScript
\texttt{export namespace} (\S\,\ref{sec:module-mapping}).

\section{Primitive Types}\label{sec:primitive-mapping}

\begin{longtable}{p{3.5cm}p{4cm}p{5cm}}
  \toprule
  \textbf{IDL Type} & \textbf{TypeScript Type} & \textbf{Notes} \\
  \midrule
  \endhead
  \texttt{boolean}                  & \texttt{boolean}    &  \\
  \texttt{char}                     & \texttt{Char}       & branded type, see \S\,\ref{sec:char-brand} \\
  \texttt{wchar}                    & \texttt{WChar}      & branded type, see \S\,\ref{sec:char-brand} \\
  \texttt{octet}                    & \texttt{number}     & 0..255 \\
  \texttt{short} / \texttt{int16}   & \texttt{number}     &  \\
  \texttt{unsigned short} / \texttt{uint16} & \texttt{number} &  \\
  \texttt{long} / \texttt{int32}    & \texttt{number}     &  \\
  \texttt{unsigned long} / \texttt{uint32}  & \texttt{number} &  \\
  \texttt{long long} / \texttt{int64} & \texttt{bigint}   & JS number is 53-bit exact \\
  \texttt{unsigned long long} / \texttt{uint64} & \texttt{bigint} & \\
  \texttt{float}                    & \texttt{number}     & IEEE 754 binary32 \\
  \texttt{double}                   & \texttt{number}     & IEEE 754 binary64 \\
  \texttt{long double}              & \texttt{LongDouble} & opaque 16-byte carrier; see \S\,\ref{sec:long-double} \\
  \texttt{string} / \texttt{wstring} & \texttt{string}    & bounded form: see \S\,\ref{sec:string-bound} \\
  \texttt{any}                      & \texttt{DdsAny}     & boxed value with type-id; see \S\,\ref{sec:any-mapping} \\
  \bottomrule
\end{longtable}

\subsection{Bounded Strings}\label{sec:string-bound}

A bounded IDL \texttt{string<N>} or \texttt{wstring<N>}
\rfckw{SHALL} produce a TypeScript \texttt{string} property as
above, plus an exported numeric constant
\texttt{$\langle$Container$\rangle$\_$\langle$Field$\rangle$\_BOUND}
emitted alongside the type declaration carrying \texttt{N}, in
the same style used for sequences
(\S\,\ref{sec:sequence-array-mapping}) and maps
(\S\,\ref{sec:map-mapping}). The bound \rfckw{SHOULD} be
checked at the runtime-validation layer; the TypeScript
\texttt{string} type cannot enforce a maximum length
structurally. Standalone bounded-string typedefs
\rfckw{SHALL} emit
\texttt{$\langle$Alias$\rangle$\_BOUND} analogously.

\subsection{Branded Character Types}\label{sec:char-brand}

IDL \texttt{char} and \texttt{wchar} hold exactly one character;
TypeScript cannot express a fixed character count structurally.
The mapping \rfckw{SHALL} therefore use branded aliases with
string-literal tags (chosen so that brands are stable across
vendor and module boundaries):

\begin{lstlisting}
export type Char  = string & { readonly __dds_brand: "char"  };
export type WChar = string & { readonly __dds_brand: "wchar" };
\end{lstlisting}

The property name \texttt{\_\_dds\_brand} and the tag values
\texttt{"char"}, \texttt{"wchar"} are normative. Construction
\rfckw{SHALL} use the factory helpers below; runtime enforcement
is the responsibility of the boundary layer.

\paragraph{\texttt{makeChar}.} The IDL \texttt{char} is an 8-bit
quantity in the range 0..255 holding a single ISO 8859-1 character
(OMG IDL 4.2 §7.2.6.2). The helper \rfckw{SHALL} accept a value if
and only if its UTF-16 code-unit length is exactly one and the
single Unicode code point lies in the inclusive range
\texttt{U+0000}..\texttt{U+00FF}; it \rfckw{SHALL} throw a
\texttt{RangeError} otherwise.

\begin{lstlisting}
export function makeChar(s: string): Char {
  if (s.length !== 1 || s.codePointAt(0)! > 0xFF) {
    throw new RangeError(
      "Char requires single ISO 8859-1 code point");
  }
  return s as Char;
}
\end{lstlisting}

\paragraph{\texttt{makeWChar}.} The IDL \texttt{wchar} is a full
Unicode scalar (OMG IDL 4.2 §7.2.6.2): the value space is
\texttt{U+0000}..\texttt{U+10FFFF} excluding the surrogate range
\texttt{U+D800}..\texttt{U+DFFF}. Because JavaScript strings are
UTF-16 code-unit sequences, scalars above the BMP occupy two
code units (a surrogate pair); length-checking \rfckw{MUST} use
code-point count, not code-unit count.

\begin{lstlisting}
export function makeWChar(s: string): WChar {
  // String iterator yields code points, not code units.
  const cps = [...s];
  if (cps.length !== 1) {
    throw new RangeError(
      "WChar requires exactly one Unicode scalar");
  }
  const cp = s.codePointAt(0)!;
  if (cp >= 0xD800 && cp <= 0xDFFF) {
    throw new RangeError(
      "WChar must not be an unpaired surrogate");
  }
  return s as WChar;
}
\end{lstlisting}

Implementations \rfckw{MUST NOT} validate \texttt{wchar} via
\texttt{string.length === 1}, which would reject every
astral-plane character.

\subsection{\texttt{long double} Handling}\label{sec:long-double}

TypeScript \texttt{number} is IEEE 754 binary64 (ECMA-262
§6.1.6.1) and cannot represent \texttt{long double} without
precision loss. IDL \texttt{long double} \rfckw{SHALL}
therefore be mapped to the opaque carrier \texttt{LongDouble}:

\begin{lstlisting}
export interface LongDouble {
  readonly __dds_brand: "long_double";
  readonly bytes:       Uint8Array;     // length === 16
}

export function makeLongDouble(bytes: Uint8Array): LongDouble {
  if (bytes.length !== 16) {
    throw new RangeError("LongDouble requires 16 bytes");
  }
  return { __dds_brand: "long_double", bytes };
}
\end{lstlisting}

The branded shape, the brand tag \texttt{"long\_double"}, and
\texttt{bytes.length === 16} are normative. The byte order
\emph{within} \texttt{bytes} is implementation-defined; wire
encoding across a DDS boundary is governed by OMG X-Types
XCDR2 §7.4.3 and is out of scope here. Arithmetic on
\texttt{LongDouble} is outside the scope of this PSM.

\subsection{\texttt{any} Mapping}\label{sec:any-mapping}

The IDL \texttt{any} type carries a value of any IDL type
together with an explicit run-time type identity (OMG IDL 4.2
§7.4.14). Because TypeScript \texttt{interface} declarations have
no runtime representation, an \emph{unwrapped} payload object
provides no basis for run-time type discovery. The code generator
therefore \rfckw{SHALL} map \texttt{any} to the boxed type
\texttt{DdsAny}, defined in Chapter~\ref{ch:descriptor-runtime}:

\begin{lstlisting}
export interface DdsAny {
  readonly typeId: string;        // qualified IDL name
  readonly value: unknown;        // the boxed payload
}
\end{lstlisting}

Producers \rfckw{SHALL} construct a \texttt{DdsAny} via
\texttt{boxAny(descriptor, value)}; consumers \rfckw{SHALL}
narrow via \texttt{unboxAny(any, descriptor)}, which checks
\texttt{any.typeId} against the descriptor name and applies the
descriptor's type-guard before returning the typed value.

\paragraph{Static unknowability.} The runtime type identity of
a \texttt{DdsAny} value is by construction not statically known
at the TypeScript surface; the corresponding
\texttt{DdsTypeRef.kind = "any"} carries no \texttt{name} or
descriptor reference. A consumer that walks a descriptor graph
and encounters an \texttt{any}-kind type-ref \rfckw{SHALL}
defer type-resolution to the boxed value's \texttt{typeId},
typically by passing it to \texttt{lookupType(any.typeId)}.
Implementations \rfckw{SHALL NOT} map \texttt{any} to TypeScript
\texttt{any} or \texttt{unknown}; either choice would defeat the
type-identity guarantee that the IDL \texttt{any} construct
requires.

\section{Struct Mapping}\label{sec:struct-mapping}

\subsection{Uniform Interface Mapping}\label{sec:struct-uniform}

An IDL \texttt{struct} declaration \rfckw{SHALL} be mapped to a
TypeScript \texttt{export interface} declaration of the same name.
Each IDL field \rfckw{SHALL} be mapped to a property of the
interface whose TypeScript type is the mapped form of the IDL
field type. This rule is independent of any annotation attached
to the struct or its fields; annotations \rfckw{SHALL} influence
the emitted Type Descriptor (\S\,\ref{sec:struct-descriptor}) and
TSDoc tags, never the shape of the interface itself.

Implementations \rfckw{SHALL NOT} promote a struct to a TypeScript
\texttt{class} on the basis of annotations; this PSM does not
emit any TypeScript class for IDL data types. The previous
``class-promotion'' design is explicitly disallowed.

\begin{specexample}
\begin{lstlisting}
struct Point { long x; long y; };
\end{lstlisting}
maps to
\begin{lstlisting}
export interface Point {
  x: number;
  y: number;
}
\end{lstlisting}
\end{specexample}

The same shape \rfckw{SHALL} be emitted for an annotated struct;
only the side-table descriptor changes.

\begin{specexample}
\begin{lstlisting}
@final
struct Point {
  @key  long x;
  @key  long y;
};
\end{lstlisting}
maps to
\begin{lstlisting}
export interface Point {
  x: number;
  y: number;
}
\end{lstlisting}
plus the descriptor of \S\,\ref{sec:struct-descriptor}.
\end{specexample}

\subsection{Type Descriptor Side-Table}\label{sec:struct-descriptor}

For every IDL \texttt{struct} the code generator \rfckw{SHALL}
additionally emit:

\begin{enumerate}
  \item A constant of type \texttt{DdsTypeDescriptor<T>}
        (Chapter~\ref{ch:descriptor-runtime}) named
        \texttt{$\langle$TypeName$\rangle$Type}, capturing the
        struct's name, extensibility kind, topic binding (if any),
        nested flag, and the ordered list of
        \texttt{DdsMemberDescriptor} entries.
  \item A type-guard function named
        \texttt{is$\langle$TypeName$\rangle$} of TypeScript type
        \texttt{(v: unknown) => v is $\langle$TypeName$\rangle$},
        whose body \rfckw{SHALL} verify the presence and type of
        every required member of \texttt{$\langle$TypeName$\rangle$}.
        Optional members \rfckw{SHALL NOT} be required by the
        guard. The guard \rfckw{SHALL NOT} validate
        \texttt{@unit}, \texttt{@min}, or \texttt{@max} bounds;
        those are TSDoc tags only.
  \item A side-effect-free registration call
        \texttt{registerType($\langle$TypeName$\rangle$Type)} that
        adds the descriptor to the global
        \texttt{DdsTypeRegistry} (used by \texttt{unboxAny} when
        the boxed \texttt{typeId} matches).
\end{enumerate}

\begin{specexample}
\begin{lstlisting}
@final
@topic("Pose")
struct Pose {
  @key  string<32>  frame_id;
  @id(1) double     x;
  @id(2) double     y;
};
\end{lstlisting}
maps to (interface plus side-table)
\begin{lstlisting}
export interface Pose {
  frame_id: string;
  x:        number;
  y:        number;
}

export const PoseType: DdsTypeDescriptor<Pose> = {
  kind:          "struct",
  name:          "Pose",
  extensibility: "final",
  nested:        false,
  topic:         "Pose",
  fields: [
    { name: "frame_id", id: 0, type: { kind: "string", bound: 32, wide: false }, key: true,  optional: false, mustUnderstand: false },
    { name: "x",        id: 1, type: { kind: "primitive", name: "double" },      key: false, optional: false, mustUnderstand: false },
    { name: "y",        id: 2, type: { kind: "primitive", name: "double" },      key: false, optional: false, mustUnderstand: false },
  ],
  typeGuard: isPose,
};

export function isPose(v: unknown): v is Pose {
  return typeof v === "object" && v !== null
      && typeof (v as any).frame_id === "string"
      && typeof (v as any).x === "number"
      && typeof (v as any).y === "number";
}

registerType(PoseType);
\end{lstlisting}
\end{specexample}

\subsection{Exception Mapping}\label{sec:struct-exception}

An IDL \texttt{exception} declaration \rfckw{SHALL} be mapped
identically to a struct
(\S\,\ref{sec:struct-uniform}–\S\,\ref{sec:struct-descriptor}),
with two additions:

\begin{enumerate}
  \item The interface declaration \rfckw{SHALL} extend a marker
        type \texttt{DdsException} (exported from
        \texttt{@zerodds/types}, structurally
        \texttt{\{ readonly \_\_dds\_exception: true \}}) so that
        TypeScript can distinguish exception payloads from
        ordinary structs at compile time.
  \item The descriptor's \texttt{kind} field \rfckw{SHALL} be
        \texttt{"exception"} rather than \texttt{"struct"}.
\end{enumerate}

For raising an IDL exception across an Operations Profile
boundary, see \S\,\ref{sec:operations-exceptions}.

\section{Union Mapping}\label{sec:union-mapping}

\subsection{Discriminated-Union Form}

An IDL \texttt{union} \rfckw{SHALL} be mapped to a TypeScript
discriminated-union type aliased by the IDL union name. Each
branch (\texttt{<case>} grammar (130) of OMG IDL 4.2) \rfckw{SHALL}
be a distinct anonymous-interface arm of the union, carrying the
discriminator value and the case-member typed property.

The discriminator property \rfckw{SHALL} be named
\texttt{discriminator}. Its TypeScript type \rfckw{SHALL} be
literal-narrowed to the case label, not to the discriminator's
broad mapped type:

\begin{itemize}[noitemsep]
  \item Discriminator type integer or \texttt{octet}: each arm's
        \texttt{discriminator} \rfckw{SHALL} be the literal numeric
        case label (e.g.\ \texttt{1}, \texttt{2}).
  \item Discriminator type \texttt{boolean}: \texttt{true} or
        \texttt{false}.
  \item Discriminator type \texttt{char}: a literal one-character
        \texttt{Char}; the helper \texttt{makeChar(\ldots)}
        \rfckw{SHOULD} be used at construction.
  \item Discriminator type \texttt{enum}: the corresponding
        \texttt{$\langle$EnumName$\rangle$.$\langle$Member$\rangle$}
        literal.
\end{itemize}

If the IDL union contains a \texttt{default} case (grammar (132)
\texttt{<default\_case>}), the implementation \rfckw{SHALL} emit
a final arm whose \texttt{discriminator} property carries the
TypeScript-mapped form of the discriminator type \emph{minus}
the union of explicit case labels:

\begin{lstlisting}
discriminator: Exclude<DiscType, ExplicitLabels>;
\end{lstlisting}

\paragraph{Compile-time exclusivity is partial.} TypeScript's
\texttt{Exclude<T, U>} narrows literal-typed unions but does
not narrow primitive types: \texttt{Exclude<number, 1\,|\,2>}
remains \texttt{number}. The compile-time exclusivity of the
default arm therefore holds for \texttt{enum} and
\texttt{boolean} discriminators (where labels are literal
subtypes of the discriminator's TypeScript type) and \emph{not}
for integer or \texttt{octet} discriminators. For the latter,
the default arm is structurally a \texttt{number} arm; runtime
discriminator-matching enforces case selection. Implementations
\rfckw{MAY} additionally tag the default arm with a brand
property (e.g.\ \texttt{readonly \_\_default:\,true}) for
narrower compile-time discrimination, but this brand is not
normative.

If the IDL union has no \texttt{default} case, the resulting
TypeScript type is closed for \texttt{enum} and \texttt{boolean}
discriminators; for integer discriminators the closure is
runtime-only.

The case labels \rfckw{SHALL} be distinct; an IDL union with two
\texttt{case}-clauses sharing a label is rejected by the IDL
parser (OMG IDL 4.2 §7.4.1.4.4.5) and never reaches the code
generator.

\begin{specexample}
\begin{lstlisting}
enum Tag { TAG_A, TAG_B };
union Variant switch (Tag) {
  case TAG_A:        long    i;
  case TAG_B:        string  s;
  default:           octet   raw;
};
\end{lstlisting}
maps to
\begin{lstlisting}
export type Variant =
  | { discriminator: Tag.TAG_A;                        i:   number }
  | { discriminator: Tag.TAG_B;                        s:   string }
  | { discriminator: Exclude<Tag, Tag.TAG_A | Tag.TAG_B>; raw: number };
\end{lstlisting}
\end{specexample}

\subsection{Multiple Labels Per Case}

A case body \rfckw{MAY} list more than one label
(\texttt{case L1: case L2:\,T m;}); the implementation
\rfckw{SHALL} expand each label into a separate arm of the
TypeScript discriminated union, each carrying the same case
member.

\subsection{Union Descriptor Side-Table}\label{sec:union-descriptor}

For every IDL \texttt{union} the code generator \rfckw{SHALL}
emit a constant of type \texttt{DdsTypeDescriptor<T>} with
\texttt{kind:\,"union"}, plus a type-guard
\texttt{is$\langle$UnionName$\rangle$}. The descriptor's
\texttt{fields} array \rfckw{SHALL} list one
\texttt{DdsMemberDescriptor} per case member; an additional
synthetic descriptor field named \texttt{discriminator} with
\texttt{id:\,0xFFFFFFFF} \rfckw{SHALL} appear first and carry
the discriminator's \texttt{DdsTypeRef}. The reserved
discriminator id (\texttt{0xFFFFFFFF}) follows the OMG XTypes
convention for synthetic member-ids and \rfckw{MUST NOT} be
used by any user-declared member; the IDL parser \rfckw{SHALL}
reject any \texttt{@id(0xFFFFFFFF)} (\S\,\ref{sec:annotation-mapping}). Each non-discriminator
descriptor entry \rfckw{SHALL} additionally carry the
\texttt{labels} property listing the matching case labels:

\begin{lstlisting}
readonly labels?: ReadonlyArray<number | string | boolean>;
\end{lstlisting}

The default case (when present) \rfckw{SHALL} have
\texttt{labels:\,undefined} and the descriptor \rfckw{SHALL}
expose a flag \texttt{hasDefault:\,true}.

\subsection{Implicit Default}

If the IDL union does not list a \texttt{default} case and the
discriminator is not exhaustively covered by the listed cases
(OMG X-Types 1.3 §7.2.2.4.4.4.5), the descriptor \rfckw{SHALL}
mark this with \texttt{hasImplicitDefault:\,true}; an
implementation \rfckw{SHALL} treat construction with an
out-of-range discriminator at the runtime layer as a
deserialisation error, signalled by the runtime, not a TypeScript
compile-time error.

\section{Enum Mapping}\label{sec:enum-mapping}

\subsection{Default Form}

An IDL \texttt{enum} \rfckw{SHALL} be mapped to a TypeScript
\texttt{const} object plus a string-literal-union type alias.
The implementation \rfckw{SHALL NOT} emit a TypeScript
\texttt{enum} declaration: TypeScript \texttt{enum} produces a
nominal type that is not assignment-compatible with plain
string literals (forcing all callers to import the enum) and
behaves inconsistently under \texttt{--isolatedModules}, which
breaks single-file transpilers (esbuild, swc).

\begin{specexample}
\begin{lstlisting}
enum Color { RED, GREEN, BLUE };
\end{lstlisting}
maps to
\begin{lstlisting}
export const Color = {
  RED:   "RED",
  GREEN: "GREEN",
  BLUE:  "BLUE",
} as const;
export type Color = (typeof Color)[keyof typeof Color];
\end{lstlisting}
\end{specexample}

The aliased type is the union of the literal-string values
(\texttt{"RED" | "GREEN" | "BLUE"}); the \texttt{const}
object provides a value-level handle for switch statements and
discriminator labels (\S\,\ref{sec:union-mapping}).

\subsection{\texttt{@value(N)}}

If a member is annotated \texttt{@value(N)}, the implementation
\rfckw{SHALL} additionally emit a numeric companion mapping
named \texttt{$\langle$EnumName$\rangle$Ordinal} so that the
numeric value is recoverable:

\begin{specexample}
\begin{lstlisting}
enum Op {
  @value(0) ADD,
  @value(1) SUB,
  @value(7) NEG
};
\end{lstlisting}
maps to
\begin{lstlisting}
export const Op = {
  ADD: "ADD", SUB: "SUB", NEG: "NEG",
} as const;
export type Op = (typeof Op)[keyof typeof Op];

export const OpOrdinal: Readonly<Record<Op, number>> = {
  ADD: 0, SUB: 1, NEG: 7,
} as const;
export const OpFromOrdinal: ReadonlyMap<number, Op> =
  new Map([[0, "ADD"], [1, "SUB"], [7, "NEG"]]);
\end{lstlisting}
\end{specexample}

If no member of the enum carries \texttt{@value(N)}, the implicit
ordinals are 0,1,2,\ldots in declaration order; the companion
\texttt{$\langle$EnumName$\rangle$Ordinal} object \rfckw{SHALL}
still be emitted with these implicit values.

\subsection{\texttt{@bit\_bound(N)}}

\texttt{@bit\_bound(N)} on an enum (XTypes 1.3 §7.2.2.4.7)
\rfckw{SHALL} be carried in the descriptor's \texttt{bitBound}
field. The TypeScript surface is unaffected; this annotation
governs the wire encoding only.

\subsection{Enum Descriptor Side-Table}\label{sec:enum-descriptor}

For every IDL \texttt{enum} the code generator \rfckw{SHALL} emit
a constant of type \texttt{DdsTypeDescriptor<$\langle$EnumName$\rangle$>}
with \texttt{kind:\,"enum"} plus the type-guard
\texttt{is$\langle$EnumName$\rangle$}. Each member \rfckw{SHALL}
appear as a \texttt{DdsMemberDescriptor} entry whose
\texttt{type.kind\,=\,"primitive"} and whose \texttt{default}
property carries the member's ordinal value. The descriptor's
\texttt{bitBound} field \rfckw{SHALL} carry the effective bit
bound (default 32).

\section{Bitset Mapping}\label{sec:bitset-mapping}

An IDL \texttt{bitset} \rfckw{SHALL} be mapped to a TypeScript
\texttt{interface} with one property per bitfield, plus
per-bitfield
\texttt{export const $\langle$Bitset$\rangle$\_$\langle$Field$\rangle$\_BITS}
declarations giving the width in bits as an integer literal. The
TypeScript property type for a bitfield of width $W$
\rfckw{SHALL} be selected as follows:

\begin{itemize}[noitemsep]
  \item $W \leq 32$: TypeScript type \texttt{number}.
  \item $W$ in $33..64$: TypeScript type \texttt{bigint}.
\end{itemize}

This rule mirrors \S\,\ref{sec:bitmask-mapping}; bit operations on
TypeScript \texttt{number} are coerced to Int32 and are unsafe for
widths above 32 bits, while \texttt{bigint} preserves
arithmetic for the full 64-bit IDL bitfield range.

\begin{specexample}
\begin{lstlisting}
bitset Flags {
  bitfield<3>  low;
  bitfield<5>  high;
  bitfield<40> wide;     // > 32 bits
};
\end{lstlisting}
maps to
\begin{lstlisting}
export interface Flags {
  low:  number;
  high: number;
  wide: bigint;
}
export const Flags_low_BITS  = 3;
export const Flags_high_BITS = 5;
export const Flags_wide_BITS = 40;
\end{lstlisting}
\end{specexample}

The width-constant \rfckw{SHALL} be a literal integer obtained by
const-evaluation of the IDL \texttt{$\langle$positive\_int\_const$\rangle$}
expression in the bitfield declaration. Implementations
\rfckw{MUST} support decimal, hexadecimal, octal, and binary
integer literal forms.

\subsection{Bitset Descriptor Side-Table}\label{sec:bitset-descriptor}

For every IDL \texttt{bitset} the code generator \rfckw{SHALL}
emit a constant of type
\texttt{DdsTypeDescriptor<$\langle$Name$\rangle$>} with
\texttt{kind:\,"bitset"}, plus a type-guard
\texttt{is$\langle$Name$\rangle$}. Each bitfield \rfckw{SHALL}
appear as a \texttt{DdsMemberDescriptor} entry whose
\texttt{type} is
\texttt{\{ kind:\,"bitfield",\,width:\,$W$ \}}. The descriptor's
\texttt{bitBound} field \rfckw{SHALL} carry the total width
(sum of bitfield widths). This PSM \rfckw{SHALL} reject bitsets
whose total width exceeds 64; the IDL parser \rfckw{SHALL}
produce a fatal error in that case.

\section{Bitmask Mapping}\label{sec:bitmask-mapping}

An IDL \texttt{bitmask} \rfckw{SHALL} be mapped to a TypeScript
\texttt{const} object whose entries are bit-shift expressions, plus
a type alias. The element-type of the constant and of the alias
\rfckw{SHALL} be selected from the effective \texttt{@bit\_bound}
of the bitmask declaration:

\begin{itemize}[noitemsep]
  \item Effective \texttt{@bit\_bound} $\leq 32$ (default 32, OMG
        IDL 4.2 §7.4.13.5): TypeScript element type
        \texttt{number}; shift expressions \rfckw{SHALL} use the
        unsigned 32-bit form
        \texttt{(1 \texttt{<<} K) \texttt{>>>} 0}.
  \item Effective \texttt{@bit\_bound} in $33..64$: TypeScript
        element type \texttt{bigint}; shift expressions
        \rfckw{SHALL} use the BigInt form \texttt{1n \texttt{<<} Kn}.
        The implementation \rfckw{MUST NOT} emit
        \texttt{number}-typed shifts for $K \geq 32$, because the
        ECMAScript \texttt{<<} operator coerces both operands to
        Int32 and would produce wrap-around values.
\end{itemize}

The two emission forms are mutually exclusive within a single
bitmask declaration; mixed-width emission is forbidden. The
declared bit width \rfckw{SHALL} be exposed as
\texttt{export const $\langle$Bitmask$\rangle$\_BIT\_BOUND}.

\begin{specexample}
\begin{lstlisting}
bitmask Permissions { READ, WRITE, EXEC };
\end{lstlisting}
maps to
\begin{lstlisting}
export const Permissions = {
  READ:  (1 << 0) >>> 0,
  WRITE: (1 << 1) >>> 0,
  EXEC:  (1 << 2) >>> 0,
} as const;
export type Permissions = number;
export const Permissions_BIT_BOUND = 32;
\end{lstlisting}
\end{specexample}

\begin{specexample}
\begin{lstlisting}
@bit_bound(48)
bitmask BigFlags {
  flag00, flag01, /* ... */ flag47
};
\end{lstlisting}
maps to
\begin{lstlisting}
export const BigFlags = {
  flag00: 1n << 0n,
  flag01: 1n << 1n,
  // ...
  flag47: 1n << 47n,
} as const;
export type BigFlags = bigint;
export const BigFlags_BIT_BOUND = 48;
\end{lstlisting}
\end{specexample}

\subsection{Bitmask Descriptor Side-Table}\label{sec:bitmask-descriptor}

For every IDL \texttt{bitmask} the code generator \rfckw{SHALL}
emit a constant of type
\texttt{DdsTypeDescriptor<$\langle$Name$\rangle$>} with
\texttt{kind:\,"bitmask"}, plus a type-guard
\texttt{is$\langle$Name$\rangle$}. Each bit value \rfckw{SHALL}
appear as a \texttt{DdsMemberDescriptor} entry whose
\texttt{default} property carries the bit position (0-based, per
OMG IDL 4.2 §7.4.13.5) and whose \texttt{type} is
\texttt{\{ kind:\,"bitfield",\,width:\,1 \}}. The descriptor's
\texttt{bitBound} field \rfckw{SHALL} carry the effective bit
bound.

\subsection{Position Annotation}

A bit value annotated with \texttt{@position(P)} (OMG IDL 4.2
§7.4.13.5) \rfckw{SHALL} use \texttt{P} as the bit position in
both the emitted shift expression and the corresponding
\texttt{DdsMemberDescriptor.default} value. Conflicting positions
(two bit values with the same explicit position) \rfckw{SHALL} be
rejected by the IDL parser.

\section{Sequence and Array Mapping}\label{sec:sequence-array-mapping}

\subsection{Bounded and Unbounded Sequences}

An IDL \texttt{sequence$\langle$T$\rangle$} (unbounded) or
\texttt{sequence$\langle$T,\,N$\rangle$} (bounded by \texttt{N})
\rfckw{SHALL} be mapped to a TypeScript
\texttt{Array$\langle$T'$\rangle$}, where \texttt{T'} is the
TypeScript-mapped form of \texttt{T} per
\S\,\ref{sec:primitive-mapping}.

The bound \texttt{N} of a bounded sequence \rfckw{SHALL} be
emitted as an exported numeric constant alongside the type
declaration, named
\texttt{$\langle$Container$\rangle$\_$\langle$Field$\rangle$\_BOUND}.
This constant carries the upper limit and \rfckw{SHOULD} be
checked at the boundary of the runtime-validation layer.

\begin{specexample}
\begin{lstlisting}
struct Sample {
  sequence<long, 10>  pixels;     // bounded
  sequence<string>    tags;       // unbounded
};
\end{lstlisting}
maps to
\begin{lstlisting}
export interface Sample {
  pixels: Array<number>;
  tags:   Array<string>;
}
export const Sample_pixels_BOUND = 10;
\end{lstlisting}
\end{specexample}

\subsection{Fixed-Size Arrays}

An IDL fixed-size array of \texttt{T} (\texttt{T name[N]})
\rfckw{SHALL} be mapped to a TypeScript
\texttt{Array$\langle$T'$\rangle$} together with an exported
numeric constant
\texttt{$\langle$Container$\rangle$\_$\langle$Field$\rangle$\_LENGTH}
that holds the array's exact length.

\begin{specexample}
\begin{lstlisting}
struct Vector3 { double v[3]; };
\end{lstlisting}
maps to
\begin{lstlisting}
export interface Vector3 { v: Array<number>; }
export const Vector3_v_LENGTH = 3;
\end{lstlisting}
\end{specexample}

\subsection{Multi-Dimensional Arrays}

A multi-dimensional IDL array of dimensions
$[N_1][N_2]\,\ldots\,[N_k]$ \rfckw{SHALL} be mapped to
nested \texttt{Array} types of depth $k$. The
length constants \rfckw{SHALL} be emitted as
\texttt{\_LENGTH\_DIM1}, \texttt{\_LENGTH\_DIM2}, \ldots\
suffixes preserving order of declaration.

\subsection{Element Mutability}

An implementation \rfckw{MAY} emit \texttt{ReadonlyArray} in lieu
of \texttt{Array} when (a) the containing struct or typedef
carries an explicit \texttt{@final} extensibility annotation
\emph{and} (b) the implementation has been configured with
\texttt{--readonly-collections}. Implementations \rfckw{SHALL
NOT} emit \texttt{ReadonlyArray} otherwise, since DDS data
semantics permit in-place mutation between samples.

\subsection{Empty Sequences}

An empty IDL sequence-value \rfckw{SHALL} round-trip through code
generation as a TypeScript empty array (\texttt{[]}). Implementations
\rfckw{MUST NOT} encode an empty sequence as \texttt{undefined} or
\texttt{null}; the property remains present with a zero-length
array.

\subsection{Maps}\label{sec:map-mapping}

An IDL \texttt{map$\langle$K,\,V$\rangle$} (unbounded) or
\texttt{map$\langle$K,\,V,\,N$\rangle$} (bounded by \texttt{N}),
grammar (199) of OMG IDL 4.2 §7.4.13.4.3.1, \rfckw{SHALL} be
mapped according to the key type \texttt{K}:

\begin{itemize}[noitemsep]
  \item If \texttt{K} maps to a TypeScript \texttt{string} type
        (including \texttt{Char}, \texttt{WChar}, branded string
        aliases), the result \rfckw{SHALL} be
        \texttt{ReadonlyMap<K',\,V'>} where \texttt{K'} and
        \texttt{V'} are the mapped forms of \texttt{K} and
        \texttt{V}. Implementations \rfckw{SHALL NOT} use a
        TypeScript \texttt{Record<\ldots>} for IDL maps because
        \texttt{Record} loses key-uniqueness validation and
        commits to enumeration order.
  \item If \texttt{K} maps to any other TypeScript type
        (\texttt{number}, \texttt{bigint}, \texttt{boolean},
        struct, etc.), the result \rfckw{SHALL} be
        \texttt{ReadonlyMap<K',\,V'>}.
\end{itemize}

The bound \texttt{N} of a bounded map \rfckw{SHALL} be emitted as
\texttt{export const $\langle$Container$\rangle$\_$\langle$Field$\rangle$\_BOUND}
analogously to bounded sequences
(\S\,\ref{sec:sequence-array-mapping}).

The descriptor entry for a map field \rfckw{SHALL} use a
\texttt{DdsTypeRef} of the new kind \texttt{"map"}:

\begin{lstlisting}
{ kind: "map"; key: DdsTypeRef; value: DdsTypeRef; bound?: number }
\end{lstlisting}

This kind \rfckw{SHALL} be added to the \texttt{DdsTypeRef} union
defined in Chapter~\ref{ch:descriptor-runtime}.

\begin{specexample}
\begin{lstlisting}
struct Index {
  map<string<64>, long, 1024>  by_name;
};
\end{lstlisting}
maps to
\begin{lstlisting}
export interface Index {
  by_name: ReadonlyMap<string, number>;
}
export const Index_by_name_BOUND = 1024;
\end{lstlisting}
\end{specexample}

An IDL key type that does not have value-equality semantics in
JavaScript (e.g.\ a struct, which is compared by reference)
\rfckw{SHALL} be flagged with the descriptor field
\texttt{keyEqualityHazard:\,true} on the corresponding
\texttt{DdsMemberDescriptor}; consumers \rfckw{SHOULD} use the
runtime helper \texttt{equalKey(a,\,b,\,keyDescriptor)} for
lookup rather than \texttt{Map.get} alone.

\paragraph{Construction.} A producer constructs the map as a
mutable \texttt{Map<K',\,V'>} and assigns it to the
\texttt{ReadonlyMap<K',\,V'>} property; this is covariantly
sound in TypeScript and is the recommended pattern. An
implementation \rfckw{MAY} provide a
\texttt{--mutable-collections} configuration flag (analogous to
\texttt{--readonly-collections} for sequences,
\S\,\ref{sec:sequence-array-mapping}) that emits
\texttt{Map<K',\,V'>} instead of \texttt{ReadonlyMap}; the flag
is non-default.

\section{Typedef Mapping}\label{sec:typedef-mapping}

\subsection{Simple Type Aliases}

An IDL \texttt{typedef T name;} \rfckw{SHALL} be mapped to a
TypeScript \texttt{export type name = T';} where \texttt{T'} is
the TypeScript-mapped form of \texttt{T}.

\begin{specexample}
\begin{lstlisting}
typedef long  Counter;
typedef string<128>  TopicName;
\end{lstlisting}
maps to
\begin{lstlisting}
export type Counter   = number;
export type TopicName = string;
\end{lstlisting}
\end{specexample}

\subsection{Sequence Aliases}

A typedef whose right-hand side is a sequence \rfckw{SHALL} produce
a TypeScript \texttt{Array}-aliased type. The bound constant of a
bounded sequence \rfckw{SHALL} be emitted as
\texttt{$\langle$Alias$\rangle$\_BOUND}.

\begin{specexample}
\begin{lstlisting}
typedef sequence<octet, 1024>  ByteBuffer;
\end{lstlisting}
maps to
\begin{lstlisting}
export type ByteBuffer = Array<number>;
export const ByteBuffer_BOUND = 1024;
\end{lstlisting}
\end{specexample}

\subsection{Array Aliases}

A typedef whose right-hand side is a fixed-size array
\rfckw{SHALL} produce a TypeScript \texttt{Array}-aliased type
with a corresponding \texttt{\_LENGTH} constant.

\subsection{Recursive Resolution}

When a typedef chain contains multiple levels
(\texttt{typedef A B; typedef B C;}), an implementation
\rfckw{SHALL} emit each declaration in source order. Each alias
\rfckw{SHALL} reference its immediate predecessor by name; cycles
\rfckw{SHALL} be rejected at IDL-parse time and \rfckw{SHALL NOT}
reach the code generator.

\subsection{Forward References}

If a typedef references a type declared later in the same
compilation unit, the implementation \rfckw{SHALL} emit all
\texttt{export type} aliases and \texttt{export interface}
declarations in any order: TypeScript resolves both interface
and type-alias names at module scope independent of source-line
order. The implementation \rfckw{SHALL NOT} emit a forward
declaration for a type alias; type aliases cannot be merged.

For \texttt{interface} declarations the same rule holds. When two
interfaces refer to each other (mutual recursion), the
implementation \rfckw{SHALL} emit them in any order without
synthesised forward stubs.

\subsection{Typedef Descriptor Side-Table}\label{sec:typedef-descriptor}

For every named IDL typedef the code generator \rfckw{SHALL}
emit a constant of type
\texttt{DdsTypeDescriptor<$\langle$AliasName$\rangle$>} with
\texttt{kind:\,"alias"} and a single
\texttt{DdsMemberDescriptor} entry named \texttt{value} whose
\texttt{type} carries the underlying \texttt{DdsTypeRef}. A
type-guard \texttt{is$\langle$AliasName$\rangle$} \rfckw{SHALL}
be emitted that delegates to the referenced type's guard.

\subsection{\texttt{@bit\_bound} on Integer Typedefs}\label{sec:bit-bound-typedef}

An integer-valued typedef carrying \texttt{@bit\_bound(N)} (XTypes
1.3 §7.2.2.4.7) \rfckw{SHALL} be mapped according to the effective
bound:

\begin{itemize}[noitemsep]
  \item $N \leq 32$: TypeScript surface unchanged
        (\texttt{number}); descriptor records \texttt{bitBound:\,N}.
  \item $N$ in $33..64$: TypeScript surface \rfckw{SHALL} switch
        to \texttt{bigint}; descriptor records \texttt{bitBound:\,N}.
\end{itemize}

\begin{specexample}
\begin{lstlisting}
@bit_bound(40) typedef long long Counter40;
\end{lstlisting}
maps to
\begin{lstlisting}
export type Counter40 = bigint;
export const Counter40Type: DdsTypeDescriptor<Counter40> = {
  kind: "alias", name: "Counter40",
  extensibility: "appendable", nested: false,
  bitBound: 40,
  fields: [{ name: "value", id: 0,
             type: { kind: "primitive", name: "int64" },
             key: false, optional: false, mustUnderstand: false }],
  typeGuard: isCounter40,
};
\end{lstlisting}
\end{specexample}

\section{Optional and Default Values}\label{sec:optional-defaults}

\subsection{Optional Fields}

A field annotated \texttt{@optional} \rfckw{SHALL} be mapped to a
TypeScript property using the TypeScript \texttt{?:} property-
suffix syntax; the property's declared type \rfckw{SHALL} be
\texttt{T' | undefined}, where \texttt{T'} is the TypeScript-mapped
form of the IDL field type. The descriptor's corresponding
\texttt{DdsMemberDescriptor} \rfckw{SHALL} have
\texttt{optional:\,true}.

\begin{specexample}
\begin{lstlisting}
@final struct Frame {
  long          seq;
  @optional
  long          retry;
};
\end{lstlisting}
maps to
\begin{lstlisting}
export interface Frame {
  seq:    number;
  retry?: number | undefined;
}

export const FrameType: DdsTypeDescriptor<Frame> = {
  kind: "struct", name: "Frame",
  extensibility: "final", nested: false,
  fields: [
    { name: "seq",   id: 0, type: { kind: "primitive", name: "int32" }, key: false, optional: false, mustUnderstand: false },
    { name: "retry", id: 1, type: { kind: "primitive", name: "int32" }, key: false, optional: true,  mustUnderstand: false },
  ],
  typeGuard: isFrame,
};
\end{lstlisting}
\end{specexample}

\subsection{Encoding of Absent Values (Informative)}

The wire encoding of an absent optional value is governed by the
encoding profile selected by the runtime, not by this
specification. The following table is provided for cross-reference
only and \rfckw{SHALL NOT} be read as a normative requirement of
this document; the authoritative source is the cited encoding
specification.

\begin{longtable}{p{3cm}p{8.5cm}}
  \toprule
  \textbf{Encoding} & \textbf{Absent-Value Form (Cross-Reference)} \\
  \midrule
  \endhead
  XCDR2          & \texttt{is\_present}=0 octet preceding the value
                   slot (OMG X-Types 1.3 §7.4.3). \\
  JSON           & The property key is omitted from the JSON object
                   (DDS-JSON profile). \\
  AMQP-Native    & The field is encoded as the AMQP \texttt{null}
                   primitive (DDS-AMQP §3.x). \\
  \bottomrule
\end{longtable}

What is \emph{normative} in the present specification is the
TypeScript-side representation: an absent optional value
\rfckw{SHALL} be modelled as the JavaScript \texttt{undefined}
literal (not \texttt{null} and not a sentinel object). This is the
value a Codegen-conformant consumer observes after a take/read.

\subsection{Default Values}

An IDL \texttt{@default(V)} annotation \rfckw{SHALL} be carried in
two places, neither of which alters the interface shape:

\begin{enumerate}
  \item The corresponding \texttt{DdsMemberDescriptor} entry's
        \texttt{default} property \rfckw{SHALL} hold the literal
        value \texttt{V}.
  \item A side-table constant
        \texttt{export const $\langle$Type$\rangle$\_$\langle$Field$\rangle$\_DEFAULT}
        \rfckw{SHALL} be emitted, typed identically to the field
        and initialised to the literal value. Consumer code that
        wishes to populate a default-bearing struct can use the
        constant directly.
\end{enumerate}

A consumer-side helper
\texttt{withDefaults<T>(descriptor:\,DdsTypeDescriptor<T>,\\
partial:\,Partial<T>):\,T} \rfckw{SHOULD} be exposed by the
descriptor runtime; it builds a fully-populated \texttt{T} by
filling absent properties from the descriptor's \texttt{default}
metadata.

\begin{specexample}
\begin{lstlisting}
struct Config {
  @default(5000)   long    timeout_ms;
  @default("on")   string  mode;
};
\end{lstlisting}
maps to
\begin{lstlisting}
export interface Config {
  timeout_ms: number;
  mode:       string;
}

export const Config_timeout_ms_DEFAULT: number = 5000;
export const Config_mode_DEFAULT:       string = "on";

export const ConfigType: DdsTypeDescriptor<Config> = {
  kind: "struct", name: "Config",
  extensibility: "appendable", nested: false,
  fields: [
    { name: "timeout_ms", id: 0, type: { kind: "primitive", name: "int32" }, key: false, optional: false, mustUnderstand: false, default: 5000 },
    { name: "mode",       id: 1, type: { kind: "string",    bound: undefined, wide: false }, key: false, optional: false, mustUnderstand: false, default: "on" },
  ],
  typeGuard: isConfig,
};
\end{lstlisting}
\end{specexample}

\section{Annotation Mapping}\label{sec:annotation-mapping}

\subsection{Mapping Surface}

This PSM does \emph{not} emit any TypeScript class decorator, and
does \emph{not} require the \texttt{experimentalDecorators} TypeScript
compiler option. Every IDL annotation is mapped to one or more of:

\begin{itemize}[noitemsep]
  \item a field of the side-table \texttt{DdsTypeDescriptor} or
        \texttt{DdsMemberDescriptor}
        (Chapter~\ref{ch:descriptor-runtime}),
  \item a side-table \texttt{export const} (e.g.\ a default value
        constant),
  \item a TSDoc tag emitted immediately above the corresponding
        declaration.
\end{itemize}

\subsection{Mapping Table}

The following table is normative. An IDL annotation that
parses but is not listed \rfckw{SHALL} cause the code generator
to emit \texttt{DDS-TS-W002} and drop the annotation; under the
\texttt{--strict-annotations} flag (which an implementation
\rfckw{MUST} provide), the same condition fires
\texttt{DDS-TS-E004} and aborts code generation instead.

\begin{longtable}{p{3.6cm}p{9.5cm}}
  \toprule
  \textbf{IDL Annotation} & \textbf{TypeScript Mapping (Descriptor / TSDoc / Const)} \\
  \midrule
  \endhead
  \texttt{@id(N)}                & \texttt{DdsMemberDescriptor.id = N}; TSDoc \texttt{@dds-id N} on the property. \\
  \texttt{@key}                  & \texttt{DdsMemberDescriptor.key = true}; TSDoc \texttt{@dds-key} on the property; runtime helper \texttt{getKey} composes the key in declaration order. \\
  \texttt{@optional}             & TypeScript property declared with \texttt{?:}\ syntax and type \texttt{T\,|\,undefined} (\S\,\ref{sec:optional-defaults}); \texttt{DdsMemberDescriptor.optional = true}. \\
  \texttt{@default(V)}           & Side-table \texttt{export const $\langle$Type$\rangle$\_$\langle$Field$\rangle$\_DEFAULT}; \texttt{DdsMemberDescriptor.default = V} (\S\,\ref{sec:optional-defaults}). \\
  \texttt{@final}                & \texttt{DdsTypeDescriptor.extensibility = "final"}. \\
  \texttt{@appendable}           & \texttt{DdsTypeDescriptor.extensibility = "appendable"} (implicit default if no extensibility annotation). \\
  \texttt{@mutable}              & \texttt{DdsTypeDescriptor.extensibility = "mutable"}; every member \rfckw{MUST} carry \texttt{@id(N)}. \\
  \texttt{@nested}               & \texttt{DdsTypeDescriptor.nested = true}; TSDoc \texttt{@dds-nested} on the type. \\
  \texttt{@topic("name")}        & \texttt{DdsTypeDescriptor.topic = "name"}; TSDoc \texttt{@dds-topic "name"}. \\
  \texttt{@must\_understand}     & \texttt{DdsMemberDescriptor.mustUnderstand = true}; TSDoc \texttt{@dds-must-understand}. \\
  \texttt{@unit("\ldots")}       & \texttt{DdsMemberDescriptor.unit = "\ldots"}; TSDoc \texttt{@dds-unit} tag. \\
  \texttt{@min(V)} / \texttt{@max(V)} & \texttt{DdsMemberDescriptor.min} / \texttt{.max = V}; TSDoc \texttt{@dds-min} / \texttt{@dds-max}. \\
  \texttt{@hashid(\ldots)}       & \texttt{DdsMemberDescriptor.hashid = \ldots}; TSDoc \texttt{@dds-hashid}. \\
  \texttt{@autoid(SEQUENTIAL|HASH)} & \texttt{DdsTypeDescriptor.autoid = "sequential" | "hash"}; member ids assigned per the descriptor rule before emission. \\
  \texttt{@verbatim(\ldots)}     & verbatim text emitted at the placement-kind position (Chapter~\ref{ch:verbatim}). \\
  \texttt{@bit\_bound(N)}        & numeric / bigint election in bitmask and bitset mappings (\S\,\ref{sec:bitmask-mapping}, \S\,\ref{sec:bitset-mapping}); recorded as \texttt{DdsTypeDescriptor.bitBound = N}. \\
  \bottomrule
\end{longtable}

\subsection{Annotation Resolution Order}

When multiple annotations apply to the same declaration, an
implementation \rfckw{SHALL} resolve them in the following order:

\begin{enumerate}
  \item Extensibility (\texttt{@final}, \texttt{@appendable},
        \texttt{@mutable}) — at most one per type.
  \item Identity (\texttt{@id}, \texttt{@autoid}) — required on
        every member of a \texttt{@mutable} type.
  \item Cardinality (\texttt{@optional}, \texttt{@key}) — applied
        per member.
  \item Validation and metadata (\texttt{@min}, \texttt{@max},
        \texttt{@unit}, \texttt{@nested}, \texttt{@must\_understand},
        \texttt{@hashid}, \texttt{@topic}) — recorded in the
        descriptor and rendered as TSDoc tags.
  \item Verbatim (\texttt{@verbatim}) — emitted last, possibly
        repeated, per Chapter~\ref{ch:verbatim}.
\end{enumerate}

A conflicting combination \rfckw{SHALL} cause the IDL parser to
reject the input with a fatal error; the code generator
\rfckw{SHALL NOT} encounter such a combination. Conflicts include
at minimum:

\begin{itemize}[noitemsep]
  \item More than one extensibility annotation
        (\texttt{@final}, \texttt{@appendable}, \texttt{@mutable})
        on a single type.
  \item Both \texttt{@autoid} on a type and \texttt{@id(N)} on
        any of its members: the explicit member id wins, but the
        parser \rfckw{SHALL} require all members to have an
        explicit \texttt{@id} or none, never a mix.
  \item Two members of a single type sharing the same explicit
        \texttt{@id(N)}.
  \item Any member declared with \texttt{@id(0xFFFFFFFF)}: this
        id is reserved for the synthetic union discriminator
        (\S\,\ref{sec:union-descriptor}).
  \item Two bit values of a single \texttt{bitmask} sharing the
        same explicit \texttt{@position(P)}.
  \item Two cases of a single \texttt{union} sharing the same
        case label.
\end{itemize}

\subsection{TSDoc Rendering}

For every annotation that has a TSDoc tag in the table above, the
implementation \rfckw{SHALL} render the tag in a TSDoc block
immediately preceding the corresponding TypeScript declaration.
This permits TypeDoc-style documentation generators to surface the
metadata without re-parsing the IDL.

\begin{specexample}
\begin{lstlisting}
struct Telemetry {
  @unit("celsius") @min(-40) @max(125)
  double temperature;
};
\end{lstlisting}
maps to
\begin{lstlisting}
export interface Telemetry {
  /**
   * @dds-unit celsius
   * @dds-min -40
   * @dds-max 125
   */
  temperature: number;
}
\end{lstlisting}
plus the descriptor entry recording \texttt{unit:\,"celsius",
min:\,-40, max:\,125}.
\end{specexample}

% ============================================================================
%  Chapter 8 — Descriptor Runtime.
% ============================================================================
\chapter{Descriptor Runtime (Normative)}\label{ch:descriptor-runtime}

This chapter defines the runtime surface that an IDL-Codegen-
Profile implementation imports from and a Descriptor-Runtime-
Profile implementation provides. No TypeScript class decorator is
part of this PSM; reflection is provided exclusively through
first-class descriptor values.

\paragraph{Runtime Package Name.} Throughout this specification
the runtime package is referenced as \texttt{@zerodds/types}.
This identifier is the ZeroDDS reference name and is
\emph{vendor-specific}; another conformant vendor \rfckw{MAY}
publish the runtime under a different scope or unscoped name. A
conformant implementation \rfckw{SHALL} export every symbol
defined in this chapter under exactly the names given (no
prefixing, no renaming, no namespacing); only the import path is
implementation-defined. Generated code emitted by a conformant
codegen \rfckw{MAY} therefore use any package path supplied via
the \texttt{--runtime-package} configuration option, defaulting to
\texttt{@zerodds/types}.

\section{Extensibility Kind}

\begin{lstlisting}
export type ExtensibilityKind = "final" | "appendable" | "mutable";
\end{lstlisting}

The three values correspond directly to OMG X-Types 1.3 §7.2.2.4.4
\texttt{ExtensibilityKind} \texttt{\{FINAL, APPENDABLE, MUTABLE\}}.

\section{Type References}\label{sec:descriptor-typeref}

A \texttt{DdsTypeRef} describes the type of a member in
side-table form, sufficient for cross-vendor TypeObject
construction.

\begin{lstlisting}
export type PrimitiveName =
  | "boolean" | "octet"
  | "int16"  | "uint16" | "int32" | "uint32"
  | "int64"  | "uint64"
  | "float"  | "double" | "longDouble"
  | "char"   | "wchar";

export type DdsTypeRef =
  | { kind: "primitive"; name: PrimitiveName }
  | { kind: "bitfield";  width: number }    // 1..64; bitset / bitmask members only
  | { kind: "string";    bound?: number; wide: boolean }
  | { kind: "sequence";  element: DdsTypeRef; bound?: number }
  | { kind: "array";     element: DdsTypeRef; lengths: ReadonlyArray<number> }
  | { kind: "map";       key: DdsTypeRef; value: DdsTypeRef; bound?: number }
  | { kind: "ref";       name: string }     // qualified IDL name; resolve via lookupType
  | { kind: "any" };
\end{lstlisting}

The \texttt{ref} kind carries the qualified IDL name of the
referenced type rather than a direct \texttt{DdsTypeDescriptor}
reference. Resolution is performed at use time via
\texttt{lookupType(name)} (\S\,\ref{sec:descriptor-reflection}).
This rule is normative and serves three goals: it permits
mutually-recursive struct/union references that cannot be
expressed in source-line-order initialisation, it makes a
\texttt{DdsTypeRef} graph JSON-roundtrippable (no cycles in the
value graph), and it aligns with the \texttt{TypeIdentifier}
form used by OMG X-Types 1.3 for type lookup.

The \texttt{bitfield} kind is used exclusively for bitset
member-references and bitmask bit-value references
(\S\,\ref{sec:bitset-mapping}, \S\,\ref{sec:bitmask-mapping});
\texttt{width} carries the bit count (1..64). Bit-bearing
members \rfckw{SHALL NOT} use \texttt{kind:\,"primitive"} since
IDL has no \texttt{intN} primitive for arbitrary widths.

\section{Type and Member Descriptors}

\begin{lstlisting}
export interface DdsMemberDescriptor {
  readonly name:           string;
  readonly id:             number;
  readonly type:           DdsTypeRef;
  readonly key:            boolean;
  readonly optional:       boolean;
  readonly mustUnderstand: boolean;
  readonly default?:       unknown;
  readonly unit?:          string;
  readonly min?:           unknown;
  readonly max?:           unknown;
  readonly hashid?:        string;
  // Union case-member only (§sec:union-descriptor).
  readonly labels?:            ReadonlyArray<number | string | boolean>;
  // Map key with non-value-equality semantics (§sec:map-mapping).
  readonly keyEqualityHazard?: boolean;
}

export type DescriptorKind =
  | "struct" | "union" | "enum" | "bitset" | "bitmask"
  | "exception" | "alias" | "interface";

export interface DdsTypeDescriptor<T = unknown> {
  readonly kind:                 DescriptorKind;
  readonly name:                 string;     // qualified IDL name
  readonly extensibility:        ExtensibilityKind;
  readonly nested:               boolean;
  readonly topic?:               string;
  readonly autoid?:              "sequential" | "hash";
  readonly bitBound?:            number;
  // Union descriptors only (see §sec:union-descriptor).
  readonly hasDefault?:          boolean;
  readonly hasImplicitDefault?:  boolean;
  readonly fields:               ReadonlyArray<DdsMemberDescriptor>;
  readonly typeGuard:            (v: unknown) => v is T;
}
\end{lstlisting}

\subsection*{Exception Marker}

\begin{lstlisting}
export interface DdsException {
  readonly __dds_exception: true;
}
\end{lstlisting}

Every IDL \texttt{exception} interface (\S\,\ref{sec:struct-exception})
\rfckw{SHALL} extend \texttt{DdsException}.

The \texttt{kind} field disambiguates the descriptor's category;
the \texttt{fields} array \rfckw{SHALL} be empty for
\texttt{enum}, \texttt{alias}, and primitive aliases. For
\texttt{interface} descriptors the operation list is carried
separately in a Service Descriptor
(\S\,\ref{sec:operations-descriptor}).

\section{Boxed \texttt{any}}\label{sec:descriptor-any}

\begin{lstlisting}
export interface DdsAny {
  readonly typeId: string;       // qualified IDL name
  readonly value:  unknown;      // the boxed payload
}
\end{lstlisting}

\section{Branded Character Helpers}

The branded-type aliases \texttt{Char} and \texttt{WChar} together
with their factory helpers \texttt{makeChar} and \texttt{makeWChar}
are normatively defined in \S\,\ref{sec:char-brand} and
\rfckw{SHALL} be exported from this package.

\section{\texttt{LongDouble}}

The opaque type \texttt{LongDouble} and its factory
\texttt{makeLongDouble} are normatively defined in
\S\,\ref{sec:long-double} and \rfckw{SHALL} be exported from this
package only when at least one consumer crate has been generated
with \texttt{--preserve-long-double}; otherwise the symbol
\rfckw{MAY} be omitted.

\section{Reflection API}\label{sec:descriptor-reflection}

\begin{lstlisting}
export function registerType(d: DdsTypeDescriptor): void;
export function lookupType(qualifiedName: string): DdsTypeDescriptor | undefined;

export function getKey<T>(
  sample: T,
  descriptor: DdsTypeDescriptor<T>,
): ReadonlyArray<unknown>;

export function getTopic<T>(
  descriptor: DdsTypeDescriptor<T>,
): string | undefined;

export function withDefaults<T>(
  descriptor: DdsTypeDescriptor<T>,
  partial: Partial<T>,
): T;

export function boxAny<T>(
  descriptor: DdsTypeDescriptor<T>,
  value: T,
): DdsAny;

export function unboxAny<T>(
  any: DdsAny,
  descriptor: DdsTypeDescriptor<T>,
): T;

export function equalKey(
  a: unknown,
  b: unknown,
  keyType: DdsTypeRef,
): boolean;

export type GuardedTypeOf<D> =
  D extends DdsTypeDescriptor<infer T> ? T : never;

export function isOneOf<DS extends ReadonlyArray<DdsTypeDescriptor>>(
  err: unknown,
  descriptors: DS,
): err is Error & { cause: GuardedTypeOf<DS[number]> };
\end{lstlisting}

\paragraph{\texttt{registerType}.} Adds the descriptor to a
process-wide registry keyed by \texttt{DdsTypeDescriptor.name}.
A re-registration with a structurally-equivalent descriptor
\rfckw{SHALL} be silently ignored; a re-registration with a
structurally-different descriptor \rfckw{SHALL} throw.

\emph{Structural equivalence} is defined as deep value-equality
of all descriptor fields except \texttt{typeGuard}, which is
compared by reference identity (functions are not deep-equals
comparable). Two codegen outputs of the same IDL input by the
same generator version \rfckw{SHALL} produce
structurally-equivalent descriptors with reference-identical
\texttt{typeGuard} functions when the runtime is shared; cross-
process or cross-bundle scenarios where two distinct
\texttt{typeGuard} closures coexist are out of scope of the
equivalence check (the registry sees only the first
registration).

\paragraph{\texttt{lookupType}.} Returns the previously-registered
descriptor for a qualified name, or \texttt{undefined}.

\paragraph{\texttt{getKey}.} Returns the values of all members
flagged \texttt{key:\,true} in declaration order. The function
\rfckw{SHALL NOT} validate the sample's shape; the caller is
expected to invoke \texttt{descriptor.typeGuard} first if they do
not trust the sample's origin.

\paragraph{\texttt{getTopic}.} Returns
\texttt{descriptor.topic} unchanged.

\paragraph{\texttt{withDefaults}.} Returns a new value of type
\texttt{T} in which any property absent from \texttt{partial} but
defined in the descriptor's member-list with a non-\texttt{undefined}
\texttt{default} \rfckw{SHALL} be filled from that default.
Properties absent from both \texttt{partial} and the descriptor
defaults remain absent.

\paragraph{\texttt{boxAny}.} Constructs a \texttt{DdsAny} with
\texttt{typeId\,=\,descriptor.name} and the supplied
\texttt{value}. The function \rfckw{SHALL} verify
\texttt{descriptor.typeGuard(value)} and throw on failure.

\paragraph{\texttt{unboxAny}.} Returns the typed value if and only
if \texttt{any.typeId\,===\,descriptor.name} and
\texttt{descriptor.typeGuard(any.value)} returns true; otherwise
throws.

\paragraph{\texttt{equalKey}.} Performs structural value
equality between two key values described by \texttt{keyType}.
For primitive, string, \texttt{bigint}, and branded character
keys equality \rfckw{SHALL} use JavaScript value-equality
(equivalent to \texttt{===}) with one normative exception:
\texttt{NaN}-valued floating-point keys \rfckw{SHALL} compare
equal to themselves (\texttt{Object.is}-style equality), and
distinct \texttt{NaN}-bit-patterns \rfckw{SHALL} compare as
equal. \texttt{+0} and \texttt{-0} \rfckw{SHALL} compare equal.
For \texttt{ref}-kind keys referring to a struct, equality
\rfckw{SHALL} be performed recursively over every member of the
referenced descriptor (overriding the JavaScript default of
reference-equality on objects). For \texttt{any}-kind keys
equality \rfckw{SHALL} compare both \texttt{typeId} and the
unboxed value via the matching descriptor. Map, sequence, and
array values are not permitted as IDL keys.

Floating-point keys \rfckw{SHOULD} be avoided in DDS topic
designs because of representation drift; this PSM does not
forbid them but documents the comparison rule for completeness.

\paragraph{\texttt{isOneOf}.} Returns \texttt{true} if and only if
\texttt{err} is an \texttt{Error} whose attached \texttt{cause}
property satisfies the \texttt{typeGuard} of any of the supplied
exception descriptors. Used in catch-handlers of generated client
methods (\S\,\ref{sec:operations-exceptions}). The implementation
\rfckw{SHALL} treat a non-\texttt{Error} \texttt{err} as
\texttt{false}, never throw.

% ============================================================================
%  Chapter 9 — Verbatim Hook.
% ============================================================================
\chapter{Verbatim Code-Generation Hook}\label{ch:verbatim}

\section{Purpose}

The X-Types 1.3 \texttt{@verbatim} annotation
(\omgref{OMG X-Types}{\S\,7.2.2.4.8}) provides an escape hatch by
which IDL authors can inject language-specific source text into
the generated output. This chapter normatively specifies how an
IDL-Codegen-Profile implementation \rfckw{SHALL} honour that
annotation when the target language is TypeScript.

\section{Annotation Form}

The X-Types annotation form is:

\begin{lstlisting}
@verbatim(language = "ts" | "*",
          placement = <PlacementKind>,
          text      = "<source-text>")
\end{lstlisting}

An implementation \rfckw{SHALL} interpret the annotation when
\texttt{language} is the literal string \texttt{"ts"} or the
wildcard \texttt{"*"}. Other values \rfckw{SHALL} cause the
annotation to be ignored at TypeScript-codegen time (other
language PSMs may handle them).

\section{Placement Kinds}

X-Types 1.3 §7.2.2.4.8 defines six placement kinds. An
IDL-Codegen-Profile implementation \rfckw{SHALL} support all six.

\begin{longtable}{p{4cm}p{8cm}}
  \toprule
  \textbf{Placement Kind} & \textbf{TypeScript Emission Position} \\
  \midrule
  \endhead
  \texttt{BEGIN\_FILE}          & first lines of the generated \texttt{.ts} file, before the leading \texttt{// Generated by \ldots}-banner. Suited for license headers. \\
  \texttt{BEFORE\_DECLARATION}  & immediately preceding the affected declaration, before any TSDoc comment. \\
  \texttt{BEGIN\_DECLARATION}   & inside the declaration body, immediately after the opening brace. Used to inject helper methods on a class declaration. \\
  \texttt{END\_DECLARATION}     & inside the declaration body, immediately before the closing brace. \\
  \texttt{AFTER\_DECLARATION}   & immediately following the affected declaration. \\
  \texttt{END\_FILE}            & last lines of the generated \texttt{.ts} file, after all declarations. Suited for default-exports or re-exports. \\
  \bottomrule
\end{longtable}

\section{Quoting and Escaping}

The verbatim \texttt{text} field is a string literal in IDL. The
implementation \rfckw{SHALL} resolve IDL-string-escape sequences
(\texttt{\textbackslash n}, \texttt{\textbackslash t},
\texttt{\textbackslash "}, \texttt{\textbackslash\textbackslash})
\emph{before} emission, so that the resulting TypeScript source
contains the unescaped text. Implementations \rfckw{SHALL NOT}
perform any TypeScript-specific re-escaping; the responsibility
for valid TypeScript syntax lies with the IDL author.

\section{Multi-Line Text}

Multi-line verbatim blocks \rfckw{SHALL} be emitted with a leading
newline and a trailing newline, so that they integrate cleanly
between adjacent declarations. Implementations \rfckw{MAY}
preserve indentation as written; they \rfckw{SHALL NOT} alter
inner whitespace.

\section{Multiple Annotations on One Declaration}

When several \texttt{@verbatim} annotations apply to the same
declaration with the same placement-kind, they \rfckw{SHALL} be
emitted in source order. The implementation \rfckw{SHALL NOT}
re-order them or coalesce them.

\section{Reference Implementation Footprint}

The reference implementation
\texttt{crates/idl-ts/src/verbatim.rs::emit\_verbatim}
exhausts all six placement kinds and is exercised by the
test-fixture set \texttt{verbatim::tests::*}.

\begin{specexample}
\begin{lstlisting}
@verbatim(language = "ts",
          placement = "BEGIN_FILE",
          text      = "// Copyright Acme Corp.\n// SPDX-License-Identifier: Apache-2.0")
struct AcmeMessage { long seq; };
\end{lstlisting}
produces
\begin{lstlisting}
// Copyright Acme Corp.
// SPDX-License-Identifier: Apache-2.0
// Generated by zerodds idl-ts. Do not edit.

export interface AcmeMessage { seq: number; }
\end{lstlisting}
\end{specexample}

% ============================================================================
%  Chapter 10 — Code Style.
% ============================================================================
\chapter{Code Style (Informative)}\label{ch:code-style}

This chapter is informative. It collects recommendations that do
not affect conformance. Normative module and constant mapping
rules are given in
Chapter~\ref{ch:type-mapping},
\S\,\ref{sec:module-mapping} and \S\,\ref{sec:constant-mapping}.

\section{Identifier Casing}

By default the code generator preserves IDL identifier casing
verbatim. A \texttt{--camelcase} configuration option may re-case
snake\_case IDL fields to camelCase TypeScript fields.

\section{Output Format}

The generated TypeScript source uses ECMAScript-Module (ESM)
syntax. A \texttt{tsconfig.json} \texttt{target} of \texttt{ES2022}
or later is recommended; the generated source compiles cleanly
under \texttt{tsc} \texttt{--strict}.

\section{File Layout}

Implementations may emit a single \texttt{.ts} file per IDL input
file (default), one file per top-level module
(\texttt{--module-per-file}, see \S\,\ref{sec:module-mapping}), or
a custom layout driven by build-tooling integrations. None of these
choices affects the conformance status of the generated code.

% ============================================================================
%  Chapter 11 — Interface and Operation Mapping (Operations Profile).
% ============================================================================
\chapter{Interface and Operation Mapping (Normative)}\label{ch:operations}

This chapter is normative within the Operations Profile
(\S\,\ref{sec:profile-operations}). Implementations that do not
claim Operations-Profile conformance \rfckw{MAY} skip emission of
the constructs defined here.

\section{Mapping Surface}

An IDL \texttt{interface} declaration carries operations
(grammar (82) \texttt{<op\_dcl>}, OMG IDL 4.2 §7.4.3.4.3.2) and
attributes ((88) \texttt{<attr\_dcl>}). Each interface
\rfckw{SHALL} be mapped to two TypeScript interface declarations
plus one Service Descriptor:

\begin{enumerate}
  \item A \emph{client} interface, suffixed
        \texttt{$\langle$IfaceName$\rangle$Client}, exposing each
        IDL operation and attribute as an asynchronous TypeScript
        method.
  \item A \emph{handler} interface, suffixed
        \texttt{$\langle$IfaceName$\rangle$Handler}, exposing the
        same set of methods. A service-side implementer
        \rfckw{SHALL} satisfy the handler interface; the framework
        delivers calls into it.
  \item A constant of type
        \texttt{ServiceDescriptor<$\langle$IfaceName$\rangle$Client,\,$\langle$IfaceName$\rangle$Handler>}
        named
        \texttt{$\langle$IfaceName$\rangle$Service},
        carrying the operation metadata
        (\S\,\ref{sec:operations-descriptor}).
\end{enumerate}

No TypeScript class is emitted for an IDL interface; both client
and handler are structural interfaces. Concrete client
construction is the responsibility of the runtime (e.g.\ a
DDS-RPC requester adapter), which produces objects satisfying the
client interface from the Service Descriptor.

Operations \rfckw{SHALL} be mapped to \texttt{Promise<T>}, not
to synchronous returns. This diverges from the OMG C++/Java/C\#
PSMs and is necessary because TypeScript runs on event-loop
runtimes where synchronous network I/O is unavailable. A
handler that completes work synchronously \rfckw{SHALL} still
expose \texttt{Promise<T>} (returning via
\texttt{Promise.resolve}).

\section{Operation Mapping}

For each IDL operation
\texttt{$\langle$op\_type$\rangle$ $\langle$name$\rangle$($\langle$params$\rangle$) [raises\,(\,$\langle$exns$\rangle$\,)];}
the generator \rfckw{SHALL} emit a method whose:

\begin{itemize}[noitemsep]
  \item return type is \texttt{Promise<R>}, where
        \texttt{R} is the TypeScript-mapped form of
        \texttt{$\langle$op\_type$\rangle$} per Chapter
        \ref{ch:type-mapping}, and \texttt{Promise<void>} for
        \texttt{void}-returning operations;
  \item parameter list is composed of the IDL \texttt{in} and
        \texttt{inout} parameters, in declaration order, each
        typed per Chapter~\ref{ch:type-mapping}; \texttt{out}
        parameters \rfckw{SHALL} be aggregated into the resolved
        return value as described below;
  \item method name is the IDL operation name preserved verbatim,
        unless the \texttt{--camelcase} option is in effect.
\end{itemize}

\paragraph{Out and InOut parameters.} If an operation declares
either \texttt{out} or \texttt{inout} parameters, the resolved
return value \rfckw{SHALL} be an object literal carrying the
declared return value as property \texttt{result} and one
property per \texttt{out}/\texttt{inout} parameter, named after
the IDL parameter. If the operation's IDL return type is
\texttt{void}, the \texttt{result} property is omitted. This rule
preserves multiple-return semantics in TypeScript, which lacks
native out-parameter passing.

\texttt{inout} parameters \rfckw{SHALL} appear in both the
method's parameter list (as input) and in the resolved return
value's object (as output). The two are independent values; the
parameter passes the initial value, the result property carries
the value updated by the handler.

\begin{specexample}
\begin{lstlisting}
interface Cursor {
  long advance(in long step, inout long position, out boolean wrapped);
};
\end{lstlisting}
maps to
\begin{lstlisting}
export interface CursorClient {
  advance(step: number, position: number):
    Promise<{ result: number; position: number; wrapped: boolean }>;
}
\end{lstlisting}
\end{specexample}

\begin{specexample}
\begin{lstlisting}
interface Calculator {
  long add(in long a, in long b) raises (Overflow);
  long divmod(in long a, in long b, out long remainder);
  oneway void reset();
};
\end{lstlisting}
maps to (selected client surface)
\begin{lstlisting}
export interface CalculatorClient {
  add(a: number, b: number): Promise<number>;
  divmod(a: number, b: number): Promise<{ result: number; remainder: number }>;
  reset(): Promise<void>;     // oneway: see below
}
\end{lstlisting}
\end{specexample}

\paragraph{Oneway operations.} An IDL \texttt{oneway} operation
\rfckw{SHALL} return \texttt{Promise<void>}. The promise
\rfckw{SHALL} resolve when the call has been handed off to the
transport, not when the handler has executed. This corresponds to
the OMG IDL 4.2 \texttt{oneway} semantics (no reply).

\paragraph{No \texttt{async} keyword in interfaces.} TypeScript
does not permit the \texttt{async} keyword on interface method
signatures (\texttt{async} is an implementation marker, not a
type marker). The generated client and handler interfaces
\rfckw{SHALL NOT} use \texttt{async} on method signatures; the
asynchronous semantics are carried entirely by the
\texttt{Promise<T>} return type. A handler implementation
(class or object satisfying the handler interface) \rfckw{MAY}
use \texttt{async} on its method bodies.

\section{Attribute Mapping}\label{sec:operations-attributes}

An IDL \texttt{attribute T name;} declaration \rfckw{SHALL} be
mapped to a getter and (unless \texttt{readonly}) a setter on
both the client and the handler interface:

\begin{lstlisting}
get_$\langle$name$\rangle$(): Promise<T'>;
set_$\langle$name$\rangle$(value: T'): Promise<void>;
\end{lstlisting}

For \texttt{readonly attribute} the setter \rfckw{SHALL NOT} be
emitted. \texttt{getraises} / \texttt{setraises} clauses
contribute to the attribute's exception set in the descriptor as
defined in \S\,\ref{sec:operations-descriptor}.

\section{Exception Mapping}\label{sec:operations-exceptions}

An IDL \texttt{exception} declaration is mapped per
\S\,\ref{sec:struct-exception}: a structural interface plus a
descriptor with \texttt{kind:\,"exception"}.

When an operation declares
\texttt{raises\,(\,E\textsubscript{1},\,\ldots,\,E\textsubscript{n}\,)},
the framework on the client side \rfckw{SHALL} reject the returned
promise with an \texttt{Error}-derived value whose attached
\texttt{cause} is one of the structural exceptions
\texttt{E\textsubscript{1}..E\textsubscript{n}}. The runtime
helper \texttt{isOneOf<E>(err, [\ldots])} \rfckw{SHALL} be
exported by \texttt{@zerodds/types} as a type guard for
narrowing in catch-handlers.

\section{Service Descriptor}\label{sec:operations-descriptor}

\begin{lstlisting}
export type ParameterMode = "in" | "out" | "inout";

export interface OperationParameterDescriptor {
  readonly name:    string;
  readonly mode:    ParameterMode;
  readonly type:    DdsTypeRef;
}

export interface OperationDescriptor {
  readonly name:        string;
  readonly oneway:      boolean;
  readonly returnType:  DdsTypeRef | { kind: "void" };
  readonly parameters:  ReadonlyArray<OperationParameterDescriptor>;
  readonly raises:      ReadonlyArray<DdsTypeDescriptor>;
}

export interface AttributeDescriptor {
  readonly name:       string;
  readonly readonly:   boolean;
  readonly type:       DdsTypeRef;
  readonly getRaises:  ReadonlyArray<DdsTypeDescriptor>;
  readonly setRaises:  ReadonlyArray<DdsTypeDescriptor>;
}

export interface ServiceDescriptor<TClient, THandler> {
  readonly name:       string;          // qualified IDL name
  readonly inherits:   ReadonlyArray<ServiceDescriptor<unknown, unknown>>;
  readonly operations: ReadonlyArray<OperationDescriptor>;
  readonly attributes: ReadonlyArray<AttributeDescriptor>;
}
\end{lstlisting}

The implementation \rfckw{SHALL} emit one
\texttt{ServiceDescriptor} constant per IDL interface, named per
the rule in \S\,\ref{sec:operations-descriptor} above. The
descriptor's \texttt{inherits} array \rfckw{SHALL} reflect the
interface's inheritance specification (grammar (78)
\texttt{<interface\_inheritance\_spec>}) in declaration order.

\section{Inheritance}\label{sec:operations-inheritance}

The TypeScript client and handler interfaces \rfckw{SHALL}
\texttt{extends} the corresponding mapped interfaces of every
parent interface listed in the IDL inheritance specification.
This propagates the operation set, attribute set, and (via
\texttt{Service\-Descriptor.inherits}) the descriptor metadata.

\begin{specexample}
\begin{lstlisting}
interface Base { long ping(); };
interface Counter : Base { long increment(); };
\end{lstlisting}
maps to
\begin{lstlisting}
export interface BaseClient {
  ping(): Promise<number>;
}
export interface BaseHandler { ping(): Promise<number>; }

export interface CounterClient extends BaseClient {
  increment(): Promise<number>;
}
export interface CounterHandler extends BaseHandler {
  increment(): Promise<number>;
}

export const CounterService:
  ServiceDescriptor<CounterClient, CounterHandler> = {
  name:       "Counter",
  inherits:   [BaseService],
  operations: [/* increment only */],
  attributes: [],
};
\end{lstlisting}
\end{specexample}

\paragraph{Multiple Inheritance.} If an IDL interface lists more
than one parent (\texttt{interface I :\,A,\,B}), the generated
client and handler \rfckw{SHALL} \texttt{extends} all parent
client / handler interfaces in declaration order:
\texttt{export interface IClient extends AClient,\,BClient \{\,\ldots\,\}}.
TypeScript inherent rules apply: a method present on more than
one parent with conflicting signatures \rfckw{SHALL} cause a
\texttt{tsc --strict} compile-time error; the IDL parser
\rfckw{SHOULD} pre-empt this by rejecting incompatible
multiple-inheritance.

The descriptor's \texttt{inherits} array \rfckw{SHALL} contain
all parent service descriptors in declaration order; no
flattening of inherited operations into the child descriptor is
performed by this PSM.

\section{Forward Declarations}

An IDL \texttt{interface} forward-declaration (grammar (75)
\texttt{<interface\_forward\_dcl>}) \rfckw{SHALL NOT} produce a
separate TypeScript declaration. Because TypeScript resolves
type names at module scope independent of source-line order, the
code generator \rfckw{SHALL} emit only the complete declarations
of every interface; declaration order in the output is
implementation-defined and \rfckw{SHALL} have no observable
semantic effect.

If the IDL input contains a forward declaration without a
matching complete declaration in the same compilation unit, the
code generator \rfckw{SHALL} emit a fatal error
(\texttt{DDS-TS-E002}, ``forward-declared interface lacks complete
declaration'').

% ============================================================================
%  Annex A — Examples (informative).
% ============================================================================
\annex{A}{informative}{IDL Examples and Generated Output}

\section*{A.1 Bounded Sequence}

\begin{lstlisting}
struct Telemetry {
  string<32>                 name;
  sequence<double, 16>        readings;
};
\end{lstlisting}
maps to
\begin{lstlisting}
export interface Telemetry {
  name:     string;
  readings: Array<number>;
}
export const Telemetry_name_BOUND     = 32;
export const Telemetry_readings_BOUND = 16;
\end{lstlisting}

\section*{A.2 Mutable Type with Explicit Ids}

\begin{lstlisting}
@mutable
struct Config {
  @id(0)  string  endpoint;
  @id(1)  long    timeout_ms;
};
\end{lstlisting}
maps to
\begin{lstlisting}
import {
  DdsTypeDescriptor, registerType,
} from "@zerodds/types";

export interface Config {
  endpoint:   string;
  timeout_ms: number;
}

export function isConfig(v: unknown): v is Config {
  return typeof v === "object" && v !== null
      && typeof (v as any).endpoint === "string"
      && typeof (v as any).timeout_ms === "number";
}

export const ConfigType: DdsTypeDescriptor<Config> = {
  kind: "struct", name: "Config",
  extensibility: "mutable", nested: false,
  fields: [
    { name: "endpoint",   id: 0,
      type: { kind: "string", wide: false },
      key: false, optional: false, mustUnderstand: false },
    { name: "timeout_ms", id: 1,
      type: { kind: "primitive", name: "int32" },
      key: false, optional: false, mustUnderstand: false },
  ],
  typeGuard: isConfig,
};

registerType(ConfigType);
\end{lstlisting}

\section*{A.3 Bitset with Width-Constants}

\begin{lstlisting}
bitset RegisterMap {
  bitfield<8>  device_id;
  bitfield<3>  status;
  bitfield<5>  reserved;
};
\end{lstlisting}
maps to
\begin{lstlisting}
export interface RegisterMap {
  device_id: number;
  status:    number;
  reserved:  number;
}
export const RegisterMap_device_id_BITS = 8;
export const RegisterMap_status_BITS    = 3;
export const RegisterMap_reserved_BITS  = 5;
\end{lstlisting}

% ============================================================================
%  Annex B — Compliance Test Suite (normative).
% ============================================================================
\annex{B}{normative}{Compliance Test Suite}

\section*{B.1 Codegen-Profile Tests}

\begin{enumerate}
  \item \textbf{B.1.1} \texttt{boolean}, \texttt{octet},
        \texttt{short}, \texttt{long}, \texttt{long long},
        \texttt{float}, \texttt{double}, \texttt{string}
        round-trip through code generation and produce TypeScript
        that compiles under \texttt{tsc --strict} without
        \texttt{experimentalDecorators}.
  \item \textbf{B.1.2} A \texttt{struct} with three fields produces
        an \texttt{export interface} with three properties, each
        typed per \S\,\ref{sec:primitive-mapping}, accompanied by
        a \texttt{$\langle$Type$\rangle$Type:\,DdsTypeDescriptor}
        constant and an \texttt{is$\langle$Type$\rangle$} guard
        per \S\,\ref{sec:struct-mapping}. No \texttt{class}
        keyword \rfckw{MUST} appear in the generator output for
        any IDL data type.
  \item \textbf{B.1.3} A \texttt{union} produces a discriminated
        union with literal-narrowed \texttt{discriminator}
        properties per case (numeric, boolean, char or enum
        labels); a \texttt{default} case maps to an
        \texttt{Exclude<\ldots>} arm; a generated
        \texttt{$\langle$Name$\rangle$Type} descriptor with
        \texttt{kind:\,"union"} carries every case label in the
        \texttt{labels} property and the \texttt{hasDefault} /
        \texttt{hasImplicitDefault} flags per
        \S\,\ref{sec:union-descriptor}.
  \item \textbf{B.1.4} An \texttt{enum} produces a TypeScript
        \texttt{const} object plus a string-literal-union type
        alias (no \texttt{enum} keyword in the output); an
        \texttt{$\langle$Enum$\rangle$Ordinal} side-table is
        emitted whether or not \texttt{@value(N)} is used; a
        descriptor with \texttt{kind:\,"enum"} and
        \texttt{bitBound:\,32} (default) is emitted per
        \S\,\ref{sec:enum-descriptor}.
  \item \textbf{B.1.5a} A \texttt{bitset} produces an
        \texttt{interface} plus per-bitfield \texttt{\_BITS}
        constants with literal integer values; bitfields of width
        $\leq 32$ use TypeScript \texttt{number}, widths $33..64$
        use \texttt{bigint}; a descriptor with
        \texttt{kind:\,"bitset"} is emitted per
        \S\,\ref{sec:bitset-descriptor}.
  \item \textbf{B.1.5b} A \texttt{bitmask} with default
        \texttt{@bit\_bound} produces a \texttt{const} object of
        unsigned 32-bit shift expressions
        (\texttt{(1\,\texttt{<<}\,K)\,\texttt{>>>}\,0}) and a
        \texttt{number} alias; with \texttt{@bit\_bound} in
        $33..64$ produces \texttt{1n\,\texttt{<<}\,Kn} expressions
        and a \texttt{bigint} alias; \texttt{@position(P)}
        annotations override implicit positions; a descriptor with
        \texttt{kind:\,"bitmask"} is emitted per
        \S\,\ref{sec:bitmask-descriptor}.
  \item \textbf{B.1.6} A \texttt{module} containing
        \texttt{typedef}, \texttt{struct}, \texttt{enum},
        \texttt{union}, \texttt{bitmask}, and \texttt{bitset}
        declarations produces a single \texttt{export namespace}
        of the same name with all members declared, per the rule
        in \S\,\ref{sec:module-mapping}. Nested modules produce
        nested namespaces; qualified IDL references resolve to
        dotted TypeScript paths.
  \item \textbf{B.1.7} Free-standing IDL constants
        (\texttt{const long}, \texttt{const double},
        \texttt{const string}, \texttt{const boolean},
        \texttt{const long long}) produce
        \texttt{export const}-declarations with the correct
        TypeScript type and a literal value, per
        \S\,\ref{sec:constant-mapping}. \texttt{long long}
        constants emit a \texttt{bigint}-suffix literal.
  \item \textbf{B.1.8} \texttt{char} and \texttt{wchar} produce
        the branded type aliases \texttt{Char} / \texttt{WChar}
        per \S\,\ref{sec:char-brand}. \texttt{makeChar} rejects
        strings whose code-unit length $\neq 1$ or whose code
        point $> \mathtt{0xFF}$. \texttt{makeWChar} accepts
        astral-plane Unicode scalars (single code point of
        length-2 UTF-16 surrogate pair) and rejects unpaired
        surrogates.
  \item \textbf{B.1.9} An IDL \texttt{long double} declaration
        produces a field of TypeScript type \texttt{LongDouble}
        per \S\,\ref{sec:long-double}; \texttt{LongDouble}
        \rfckw{MUST} have exactly the normative shape
        \texttt{\{ readonly \_\_dds\_brand:\,"long\_double";
        readonly bytes:\,Uint8Array \}} with
        \texttt{bytes.length === 16}. The byte order inside
        \texttt{bytes} is implementation-defined and \rfckw{SHALL
        NOT} be tested by this conformance suite. The generator
        \rfckw{MUST NOT} substitute \texttt{number} and
        \rfckw{MUST NOT} abort code generation on this construct.
  \item \textbf{B.1.10} An IDL \texttt{any} declaration produces
        a property of TypeScript type \texttt{DdsAny} per
        \S\,\ref{sec:any-mapping}; type \texttt{any} or
        \texttt{unknown} \rfckw{MUST NOT} appear in the generator
        output for an IDL \texttt{any} member.
  \item \textbf{B.1.11} An IDL \texttt{struct} carrying
        \texttt{@key}, \texttt{@id}, \texttt{@final},
        \texttt{@appendable}, \texttt{@mutable},
        \texttt{@optional}, \texttt{@default}, or \texttt{@topic}
        produces the same shape of \texttt{export interface} as
        an un-annotated struct; annotation values surface only in
        the side-table descriptor and (where applicable) TSDoc
        tags. No TypeScript \texttt{class} is emitted for any IDL
        data type.
  \item \textbf{B.1.12} A free-standing IDL \texttt{exception}
        declaration produces an \texttt{export interface}
        extending \texttt{DdsException} and a descriptor with
        \texttt{kind:\,"exception"} per
        \S\,\ref{sec:struct-exception}.
  \item \textbf{B.1.13} A named IDL \texttt{typedef} produces a
        TypeScript \texttt{export type} alias and a descriptor
        with \texttt{kind:\,"alias"} per
        \S\,\ref{sec:typedef-descriptor}; an integer typedef
        carrying \texttt{@bit\_bound(N)} with $N \in 33..64$
        switches the alias to \texttt{bigint} per
        \S\,\ref{sec:bit-bound-typedef}.
  \item \textbf{B.1.14} An IDL \texttt{module} containing a
        free-standing \texttt{const}, a \texttt{struct}, a
        \texttt{union} with \texttt{default}, an \texttt{enum}
        with \texttt{@value(N)}, a \texttt{bitset} mixing widths
        $\leq 32$ and $> 32$, a \texttt{bitmask} with
        \texttt{@bit\_bound(40)}, and a \texttt{typedef}
        compiles cleanly under \texttt{tsc --strict} without
        \texttt{experimentalDecorators}.
  \item \textbf{B.1.15} An unrecognised but syntactically-valid
        IDL annotation causes the generator to emit
        \texttt{DDS-TS-W002} and to drop the annotation; code
        generation \rfckw{MUST NOT} abort.
  \item \textbf{B.1.16} An IDL \texttt{map<K, V>} field produces
        a TypeScript \texttt{ReadonlyMap<K', V'>} property; a
        bounded map produces an additional
        \texttt{$\langle$Container$\rangle$\_$\langle$Field$\rangle$\_BOUND}
        constant; the descriptor entry uses
        \texttt{type.kind:\,"map"} per
        \S\,\ref{sec:map-mapping}. A struct-typed key sets
        \texttt{keyEqualityHazard:\,true}.
\end{enumerate}

\section*{B.2 Descriptor-Runtime Tests}

\begin{enumerate}
  \item \textbf{B.2.1} The package \texttt{@zerodds/types} exports
        the descriptor types and helper functions defined in
        Chapter~\ref{ch:descriptor-runtime}:
        \texttt{DdsTypeDescriptor}, \texttt{DdsMemberDescriptor},
        \texttt{DdsTypeRef}, \texttt{DdsAny},
        \texttt{DdsException}, \texttt{DescriptorKind},
        \texttt{ExtensibilityKind}, \texttt{PrimitiveName},
        \texttt{Char}, \texttt{WChar}, \texttt{makeChar},
        \texttt{makeWChar}, \texttt{registerType},
        \texttt{lookupType}, \texttt{getKey}, \texttt{getTopic},
        \texttt{withDefaults}, \texttt{boxAny},
        \texttt{unboxAny}, \texttt{equalKey}, \texttt{isOneOf}.
  \item \textbf{B.2.2} \texttt{lookupType("FQN")} returns the
        descriptor previously registered with that name, or
        \texttt{undefined}.
  \item \textbf{B.2.3} \texttt{getKey(sample, descriptor)}
        returns an array containing the values of all
        \texttt{key:\,true} members in declaration order; for a
        struct without key members it returns \texttt{[]}.
  \item \textbf{B.2.4} \texttt{boxAny(descriptor, value)} throws
        when \texttt{descriptor.typeGuard(value)} returns false;
        otherwise it returns
        \texttt{\{ typeId:\,descriptor.name, value\,\}}.
  \item \textbf{B.2.5} \texttt{unboxAny(any, descriptor)} returns
        the typed value when both \texttt{any.typeId} and
        \texttt{descriptor.typeGuard(any.value)} match, and
        throws otherwise.
  \item \textbf{B.2.6} \texttt{withDefaults(descriptor,\,partial)}
        fills absent properties from descriptor-recorded
        defaults; properties with no default remain absent.
  \item \textbf{B.2.7} \texttt{makeWChar("\textbackslash u\{1F600\}")}
        (the smiling-face emoji, two UTF-16 code units, one Unicode
        scalar) succeeds; \texttt{makeChar(\textquotedbl{}ß\textquotedbl{})}
        succeeds (\texttt{U+00DF}, in ISO 8859-1);
        \texttt{makeChar(\textquotedbl{}€\textquotedbl{})} throws
        (\texttt{U+20AC}, outside ISO 8859-1).
  \item \textbf{B.2.8} \texttt{equalKey} returns true for two
        struct-typed keys with deeply-equal members; returns false
        for two reference-equal but structurally-different keys
        only if the members differ.
  \item \textbf{B.2.9} \texttt{isOneOf(err, [E1Type, E2Type])}
        returns true for a rejection produced by an Operations-
        Profile client when the underlying IDL operation raised
        \texttt{E1} or \texttt{E2}; returns false otherwise; never
        throws on non-\texttt{Error} input.
\end{enumerate}

\section*{B.3 Operations-Profile Tests}

\begin{enumerate}
  \item \textbf{B.3.1} An IDL \texttt{interface I} produces both
        \texttt{IClient} and \texttt{IHandler} TypeScript
        \texttt{interface} declarations and an \texttt{IService:
        ServiceDescriptor} constant. No \texttt{class} keyword
        \rfckw{MUST} appear in the generator output for any IDL
        \texttt{interface}.
  \item \textbf{B.3.2} An operation
        \texttt{R foo(in T1 a, out T2 b);} produces a method
        \texttt{foo(a:\,T1):\,Promise<\{ result:\,R; b:\,T2 \}>}
        on both client and handler.
  \item \textbf{B.3.3} An operation
        \texttt{void bar(in T x) raises (E);} produces a method
        \texttt{bar(x:\,T):\,Promise<void>} and the descriptor's
        \texttt{raises} list contains the \texttt{E} descriptor.
  \item \textbf{B.3.4} A \texttt{oneway void baz();} operation
        produces a method \texttt{baz():\,Promise<void>} and the
        descriptor's \texttt{oneway} flag is true.
  \item \textbf{B.3.5} A \texttt{readonly attribute T name;}
        produces \texttt{get\_name():\,Promise<T>} but no
        \texttt{set\_name}.
  \item \textbf{B.3.6} An interface inheritance chain
        \texttt{interface B :\,A \{\}} produces
        \texttt{BClient extends AClient},
        \texttt{BHandler extends AHandler}, and
        \texttt{BService.inherits[0]\,===\,AService}.
  \item \textbf{B.3.7} Multiple inheritance
        \texttt{interface I :\,A,\,B \{\}} produces
        \texttt{IClient extends AClient,\,BClient} and
        \texttt{IService.inherits} of length 2 in declaration
        order. A signature conflict between \texttt{A} and
        \texttt{B} on the same operation name \rfckw{SHALL} cause
        \texttt{tsc --strict} to fail; the IDL parser
        \rfckw{SHOULD} pre-empt this with a parser-level error.
  \item \textbf{B.3.8} A forward-only interface declaration
        without a matching complete declaration in the same
        compilation unit causes \texttt{DDS-TS-E002}.
  \item \textbf{B.3.9} A \texttt{readonly attribute T name} with
        \texttt{getraises\,(\,E\,)} produces a method
        \texttt{get\_name():\,Promise<T>} and the descriptor's
        \texttt{attributes[i].getRaises} contains the descriptor
        of \texttt{E}.
\end{enumerate}

\section*{B.4 Reference Test-Harness}

The reference test-harness is hosted at
\url{https://github.com/zero-objects/zero-dds/}
under \texttt{crates/idl-ts/} (Rust unit tests against the codegen
output) and verifies B.1, B.2, and (when the
Operations Profile is claimed) B.3 by parsing the generated
TypeScript with \texttt{tsc} \texttt{--noEmit --strict} and
exercising the descriptor-runtime helpers in a Node.js test
runner.

% ============================================================================
%  Annex C — WASM-Bindings Profile (normative).
% ============================================================================
\annex{C}{normative}{WASM-Bindings Profile}\label{sec:wasm-reservation}

This annex normatively specifies the WASM-Bindings Profile of
DDS-TS. An implementation that claims conformance to the Profile
\rfckw{SHALL} satisfy every requirement in this annex.

\section*{C.1 Handle and Sample Types}\label{sec:wasm-types}

The handle and sample types referenced throughout this annex are
normatively defined as follows. They \rfckw{SHALL} be exported
from the implementation's binding module under the names below.

\subsection*{C.1.1 Branded Handle Types}

The handle types \rfckw{SHALL} be defined using string-literal
brand tags rather than module-local \texttt{unique symbol}
declarations. String-literal brand tags compare structurally
across module and vendor boundaries; \texttt{unique symbol}-
based brands do not, and would prevent a handle produced by
vendor A's binding from satisfying vendor B's handle parameter
type even when the wire-level resource is identical.

\begin{lstlisting}
export type ParticipantHandle = number & { readonly __dds_brand: "participant" };
export type TopicHandle       = number & { readonly __dds_brand: "topic"       };
export type PublisherHandle   = number & { readonly __dds_brand: "publisher"   };
export type SubscriberHandle  = number & { readonly __dds_brand: "subscriber"  };
export type DataWriterHandle  = number & { readonly __dds_brand: "writer"      };
export type DataReaderHandle  = number & { readonly __dds_brand: "reader"      };
\end{lstlisting}

The brand-tag values (\texttt{"participant"}, \texttt{"topic"},
\texttt{"publisher"}, \texttt{"subscriber"}, \texttt{"writer"},
\texttt{"reader"}) are normative; conformant implementations
\rfckw{SHALL NOT} substitute alternative tag values. The
property-name \texttt{\_\_dds\_brand} is also normative. This
guarantees that two conformant TypeScript bindings — even from
distinct vendors and distinct npm packages — produce
type-compatible handle types.

Each handle is a non-negative 32-bit integer index into the
host's resource table. Branding prevents accidental
cross-assignment between handle kinds at compile time; the
brand is a TypeScript-only marker without runtime effect. A
handle value of \texttt{0} is reserved as the ``invalid
handle'' sentinel — this is a \emph{runtime-enforced}
invariant: the type system cannot prevent
\texttt{0 as DataReaderHandle} from compiling, so any operation
receiving \texttt{0} \rfckw{SHALL} reject with an error at
runtime.

\subsection*{C.1.2 Sample Type}

\begin{lstlisting}
/** OMG DDSI-RTPS GUID (16 bytes = 12-byte GuidPrefix + 4-byte EntityId). */
export interface DdsGuid {
  readonly prefix:   Uint8Array;   // length === 12
  readonly entityId: number;       // 32-bit unsigned
}

export function makeDdsGuid(prefix: Uint8Array, entityId: number): DdsGuid {
  if (prefix.length !== 12) {
    throw new RangeError("DdsGuid.prefix requires 12 bytes");
  }
  if (!Number.isInteger(entityId) || entityId < 0 || entityId > 0xFFFFFFFF) {
    throw new RangeError("DdsGuid.entityId requires uint32");
  }
  return { prefix, entityId };
}

export interface SampleInfo {
  readonly validData:         boolean;
  readonly sampleState:       "read" | "not_read";
  readonly viewState:         "new"  | "not_new";
  readonly instanceState:     "alive" | "not_alive_disposed" | "not_alive_no_writers";
  readonly sourceTimestampNs: bigint;
  readonly sequenceNumber:    bigint;
  readonly publicationHandle: DdsGuid;
  readonly instanceHandle:    bigint;
}

export interface Sample {
  readonly bytes: Uint8Array;            // XCDR2-encoded payload
  readonly info:  SampleInfo;
}
\end{lstlisting}

The \texttt{bytes} field carries the XCDR2 octet sequence
exactly as it crosses the wire. The \texttt{info} object is a
strict subset of the OMG DDS 1.4 \texttt{SampleInfo} structure
(§2.2.2.5.4); fields not enumerated above are not exposed by
this Profile. \texttt{publicationHandle} carries the full
DDSI-RTPS GUID of the originating writer (12-byte
\texttt{GuidPrefix} carrying vendor-id, host-id, app-id, and
instance-id components per OMG DDSI-RTPS 2.5 §8.2.4, plus
4-byte \texttt{EntityId}); truncation to a 64-bit subset
\rfckw{SHALL NOT} be performed because it loses the
vendor/host components and risks hash collisions in
multi-vendor or multi-host deployments. \texttt{instanceHandle}
is opaque and remains a 64-bit handle.

\subsection*{C.1.3 Listener Callback Type}

\begin{lstlisting}
export type DataAvailableCallback = (r: DataReaderHandle) => void;
\end{lstlisting}

\section*{C.2 Required Operations}

The WASM-Bindings Profile \rfckw{SHALL} expose the following
DDS DCPS data-path operations as TypeScript-callable functions.
The signatures are given in TypeScript and are normative; the
underlying WebAssembly export-name is implementation-defined.

\subsection*{C.2.1 Participant and Topic}

\begin{lstlisting}
function createParticipant(domainId: number): ParticipantHandle;
function deleteParticipant(p: ParticipantHandle): void;

function createTopic(p: ParticipantHandle,
                     name: string,
                     typeName: string): TopicHandle;
function deleteTopic(t: TopicHandle): void;
\end{lstlisting}

\subsection*{C.2.2 Publisher and Subscriber}

\begin{lstlisting}
function createPublisher (p: ParticipantHandle): PublisherHandle;
function createSubscriber(p: ParticipantHandle): SubscriberHandle;
function deletePublisher (pub: PublisherHandle): void;
function deleteSubscriber(sub: SubscriberHandle): void;
\end{lstlisting}

\subsection*{C.2.3 DataWriter and DataReader}

\begin{lstlisting}
function createDataWriter(pub: PublisherHandle,
                          topic: TopicHandle): DataWriterHandle;
function createDataReader(sub: SubscriberHandle,
                          topic: TopicHandle): DataReaderHandle;

function writeSample (w: DataWriterHandle, xcdr2: Uint8Array): void;
function takeSamples (r: DataReaderHandle, max: number): ReadonlyArray<Sample>;

function deleteDataWriter(w: DataWriterHandle): void;
function deleteDataReader(r: DataReaderHandle): void;
\end{lstlisting}

\subsection*{C.2.4 Listener Registration}

\begin{lstlisting}
function setDataAvailableListener(
  r: DataReaderHandle,
  cb: DataAvailableCallback | null,
): void;
\end{lstlisting}

A \texttt{cb} value of \texttt{null} \rfckw{SHALL} unregister any
previously-installed listener. Listener invocation \rfckw{SHALL}
occur on the host event loop, never inside a synchronous
\texttt{writeSample} call.

\section*{C.3 Wire-Format Boundary}

\subsection*{C.3.1 XCDR2 Sample Bytes}

Sample data \rfckw{SHALL} cross the WASM boundary as
\texttt{Uint8Array} of XCDR2-serialised octets, encoded according
to OMG X-Types 1.3 §7.4.3. The selected encapsulation kind
(\texttt{CDR2\_BE} or \texttt{CDR2\_LE}) \rfckw{SHALL} be
configured via the codegen flag \texttt{--xcdr-endianness=BE|LE},
defaulting to \texttt{LE} (the wire-form predominant in
practice). Implementations
\rfckw{MUST NOT} expose JavaScript-native object representations
as the canonical wire form; the binding layer is purely a
WebAssembly-glue layer.

\subsection*{C.3.2 Memory Ownership}

The TypeScript caller \rfckw{SHALL} retain ownership of any
\texttt{Uint8Array} passed across the boundary; the WASM module
\rfckw{SHALL} treat such buffers as borrowed and \rfckw{SHALL NOT}
retain references after the call returns. For samples returned
from \texttt{takeSamples}, the WASM module \rfckw{SHALL} allocate
fresh \texttt{Uint8Array} buffers owned by JavaScript.

\subsection*{C.3.3 \texttt{DdsAny} Across the Boundary}

A topic-type that contains an IDL \texttt{any} field
(\S\,\ref{sec:any-mapping}) crosses the WASM boundary as a fully
serialised XCDR2 payload exactly like any other topic-type;
\texttt{DdsAny} is a TypeScript-side helper for narrowing on
take and for boxing on write, not a separate wire object. The
producer \rfckw{SHALL} call \texttt{boxAny(\ldots)} on the
TypeScript side, then serialise the resulting compound topic
through the codegen-emitted serialiser; the consumer \rfckw{SHALL}
deserialise the topic, then call \texttt{unboxAny(\ldots)} for
narrowing.

\section*{C.4 Browser vs Node.js Bindings}

\subsection*{C.4.1 Browser Profile}

A browser-side conformant implementation \rfckw{SHALL} expose the
operations of \S\,C.2 via a single \texttt{.wasm} module compiled
with \texttt{wasm-bindgen} or an equivalent toolchain. The
underlying transport \rfckw{SHALL} use WebTransport, WebSocket,
or HTTP-3 for off-host communication; native UDP-multicast is
not available in browsers and \rfckw{MUST NOT} be required.

\subsection*{C.4.2 Node.js Profile}

A Node.js-side conformant implementation \rfckw{MAY} additionally
provide a native \texttt{.node} module via N-API
(\texttt{napi-rs}) for performance-sensitive deployments. The
exposed API \rfckw{SHALL} match \S\,C.2 signature-for-signature,
so a TypeScript application can be retargeted between browser
and Node.js without source-code changes outside the import
statement.

\section*{C.5 Conformance Tests}

\subsection*{C.5.1 Round-Trip Publish/Subscribe}

A WASM-Bindings-Profile implementation \rfckw{SHALL} pass the
following round-trip test:

\begin{enumerate}
  \item Create a Participant in domain~0, a Topic \texttt{"Echo"}
        with type-name \texttt{"PrimitiveSample"}, a DataWriter,
        and a DataReader on the Topic.
  \item Encode a sample of an IDL struct
        \texttt{struct PrimitiveSample \{ long v;\,\};} via the
        IDL-Codegen-Profile output and write the resulting
        XCDR2 bytes via \texttt{writeSample}.
  \item Within 1\,s, \texttt{takeSamples} \rfckw{SHALL} return an
        array of length 1 whose \texttt{info.validData} is
        \texttt{true} and whose decoded payload \texttt{v} equals
        the published value.
\end{enumerate}

\subsection*{C.5.2 Listener Callback}

After registering a \texttt{DataAvailable} listener, writing a
sample \rfckw{SHALL} cause the listener to be invoked at least
once. Listener delivery \rfckw{SHALL NOT} occur synchronously
within \texttt{writeSample}.

\subsection*{C.5.3 Resource Cleanup}

Calling \texttt{deleteParticipant} \rfckw{SHALL} release all
WebAssembly memory associated with the Participant; subsequent
calls referencing any of its derived handles \rfckw{SHALL}
throw a \texttt{RangeError} naming the invalidated handle. A
double-delete of any handle \rfckw{SHALL} also throw, never
silently succeed.

\subsection*{C.5.4 \texttt{DdsAny} Round-Trip}

A topic containing a field of IDL type \texttt{any} \rfckw{SHALL}
round-trip via the producer/consumer pattern in \S\,C.3.3 such
that the consumer's \texttt{unboxAny} call produces a value
equal to the original under
\texttt{descriptor.typeGuard}.

\section*{C.6 Reservation of Future Extensions}

The following identifiers are reserved for future revisions of
this Profile and \rfckw{SHALL NOT} be re-used by Beta
implementations for unrelated purposes:
\texttt{createGuardCondition},
\texttt{createWaitSet}, \texttt{createQueryCondition},
\texttt{getInstance}.

The identifier \texttt{registerType} is \emph{not} reserved here:
it is used by the Descriptor-Runtime Profile
(\S\,\ref{sec:profile-runtime}) for its specified semantics
(\S\,\ref{sec:descriptor-reflection}).

% ============================================================================
%  Annex D — Diagnostic Codes (normative).
% ============================================================================
\annex{D}{normative}{Diagnostic Codes}\label{annex:diagnostics}

A conformant code generator \rfckw{SHALL} use the following
diagnostic codes when emitting the conditions enumerated below.
The code, severity, and trigger condition are normative; the
human-readable message text is implementation-defined and
\rfckw{SHOULD} cite the source location of the offending IDL
construct.

\begin{longtable}{p{2.4cm}p{1.6cm}p{9.5cm}}
  \toprule
  \textbf{Code} & \textbf{Severity} & \textbf{Trigger} \\
  \midrule
  \endhead
  \texttt{DDS-TS-E001} & fatal &
    \emph{Reserved} for future use; conformant implementations
    \rfckw{SHALL NOT} re-use this code. \\
  \texttt{DDS-TS-E002} & fatal &
    A forward-declared interface
    (\texttt{<interface\_forward\_dcl>}) lacks a matching complete
    declaration in the same compilation unit
    (\S\,\ref{ch:operations}, Forward Declarations). \\
  \texttt{DDS-TS-E003} & fatal &
    The IDL parser encountered a conflicting annotation
    combination as enumerated in
    \S\,\ref{sec:annotation-mapping}; the parser \rfckw{SHALL}
    reject the input. \\
  \texttt{DDS-TS-E004} & fatal &
    Under \texttt{--strict-annotations}, an unrecognised IDL
    annotation was encountered (else fires as
    \texttt{DDS-TS-W002}, \S\,\ref{sec:annotation-mapping}). \\
  \texttt{DDS-TS-W001} & warning &
    \emph{Reserved} for future use; conformant implementations
    \rfckw{SHALL NOT} re-use this code. \\
  \texttt{DDS-TS-W002} & warning &
    A syntactically-valid IDL annotation not listed in
    \S\,\ref{sec:annotation-mapping} was encountered; the
    annotation is dropped, code generation continues. \\
  \texttt{DDS-TS-W003} & warning &
    An IDL \texttt{map} key type with non-value-equality
    JavaScript semantics (\S\,\ref{sec:map-mapping}); descriptor
    entry carries \texttt{keyEqualityHazard:\,true}. \\
  \texttt{DDS-TS-W004} & warning &
    An IDL \texttt{union} with no \texttt{default} and
    non-exhaustive cases (\S\,\ref{sec:union-mapping}); descriptor
    carries \texttt{hasImplicitDefault:\,true}. \\
  \texttt{DDS-TS-I001} & info &
    Use of IDL \texttt{long double}; mapped to the opaque
    \texttt{LongDouble} carrier (\S\,\ref{sec:long-double}). \\
  \bottomrule
\end{longtable}

Codes \texttt{E001..E099}, \texttt{W001..W099}, and
\texttt{I001..I099} are reserved for this specification; codes
\texttt{E100+}, \texttt{W100+}, \texttt{I100+} are available for
implementation-specific diagnostics. A diagnostic \rfckw{SHALL}
carry at minimum the code, severity, file path, source line and
column, and message text; the surface format
(terminal/LSP/JSON) is implementation-defined.

\end{document}
